\documentclass[a4paper,12pt]{article}
\usepackage[utf8]{inputenc}
\usepackage[german]{babel}
\usepackage{amsmath}
\usepackage{amssymb}
\usepackage{amsthm}
\usepackage{geometry}
\usepackage{graphicx}
\usepackage[hidelinks]{hyperref}
\usepackage{listings}
\usepackage{xcolor}
\usepackage{enumitem}
\usepackage{booktabs}
\usepackage{array}
\usepackage{caption}
\usepackage[protrusion=true,expansion=false,final]{microtype}

\geometry{margin=2.5cm}

% Verhindert, dass Text über den rechten Rand hinausragt
\setlength{\emergencystretch}{3em}
\tolerance=1500
\hbadness=1500

\theoremstyle{definition}
\newtheorem{definition}{Definition}
\newtheorem{beispiel}{Beispiel}
\newtheorem{satz}{Satz}
\newtheorem{lemma}{Lemma}
\newtheorem{bemerkung}{Bemerkung}
\newtheorem{korollar}{Korollar}
\newtheorem{hypothese}{Hypothese}
\renewcommand{\proofname}{Beweis}
\newenvironment{beweis}{\begin{proof}}{\end{proof}}

\setlist[itemize,1]{label=--}

% ---------- Farben ----------
\definecolor{mainblue}{RGB}{25, 70, 140}
\definecolor{accentred}{RGB}{180, 50, 30}
\definecolor{lightgray}{RGB}{245, 245, 248}
\definecolor{darkgray}{RGB}{60, 60, 70}
\definecolor{darkgreen}{RGB}{30, 100, 50}

\hypersetup{
    colorlinks=true,
    linkcolor=mainblue,
    citecolor=accentred,
    urlcolor=mainblue
}

\title{Zur Belastbarkeit von Softwarevertr\"agen:\\
Eine formale Analyse der Unsicherheit, des Risikos\\
und der Qualifikationsanforderungen in der\\
Softwareentwicklung}
\author{Stephan Epp\\
\small Universit\"at Bielefeld\\
\small \texttt{Stephan\_Epp@web.de}}
\date{4. Juni 2026}

\begin{document}

\maketitle

\begin{abstract}
Softwaresysteme unterscheiden sich grundlegend von mathematischen Systemen: Sie sind in ihrer
Gesamtheit weder vollst\"andig formal beschreibbar noch vollst\"andig verifikationsf\"ahig.
Diese Arbeit untersucht, was daraus f\"ur die Belastbarkeit von Vertr\"agen zwischen Kunden
und Softwareentwicklern folgt. Wir f\"uhren ein formales Modell der \emph{Vertragsbelastbarkeit}
$B(t, u, r)$ ein, definieren Klassen von Softwarevertr\"agen nach Unsicherheits- und
Komplexit\"atsgrad und beweisen formal, dass die Belastbarkeit mit wachsender Projektlaufzeit,
steigender Benutzerzahl und zunehmenden Anforderungen monoton abnimmt. Ein zentrales Ergebnis
dieser Arbeit ist die formale Begr\"undung des Satzes \emph{„Software hei\ss{}t Risiko"}:
Kein noch so umfangreiches Testen vor der Inbetriebnahme kann absolute Sicherheit f\"ur den
Produktionseinsatz garantieren. Dieses inhärente Restrisiko muss von allen Beteiligten
bewusst eingeplant werden. Dar\"uber hinaus wird die Frage gestellt, welche
Berufsqualifikationen f\"ur die Beurteilung von Softwarevertr\"agen ben\"otigt werden, und
es wird ein multidisziplin\"ares Kompetenzmodell vorgeschlagen. Alle Ergebnisse werden
formal bewiesen und durch Matplotlib-Visualisierungen veranschaulicht.
\end{abstract}

\tableofcontents
\newpage

% ===========================================================================
\section{Einf\"uhrung}
% ===========================================================================

Mathematische Systeme sind mathematisch, beschreibbar und l\"osbar und viele Aussagen zu
diesen Systemen lassen sich formal beweisen und machen das System damit sicher oder auch
vorhersagbar. Denn die Antwort bei gewissen Anregungen kann berechnet oder abgeleitet werden.

F\"ur Softwaresysteme trifft das nicht zu. Softwaresysteme lassen sich nicht in allen
Bereichen formal beschreiben und beweisen oder formal verifizieren. Das hat zur Folge, dass
nicht auf jede Anregung hin die Antwort des Systems klar ist oder vorhersagbar ist.

Die zweite Herausforderung besteht darin, dass oft die Kunden, f\"ur die das Softwaresystem
entwickelt wird, selber noch gar nicht alle Anforderungen kennen, die das Softwaresystem
befriedigen muss.

Die dritte und vierte Herausforderung besteht darin, dass auch die physikalische Umgebung
oder auch die Benutzer, die das System sp\"ater einsetzen, nicht in Randbereichen und nicht
vollst\"andig bekannt sind.

Betrachtet man diese vier Herausforderungen f\"ur die Entwicklung von Softwaresystemen,
stellt sich die Frage, wie zuverl\"assig und belastbar Vertr\"age zwischen Kunden und
Softwareingenieuren wirklich sein k\"onnen.

Gibt es ein Ma\ss{} f\"ur die Belastbarkeit eines Vertrages zwischen einem Kunden und
Softwareentwicklern, die die Software f\"ur den Kunden entwickeln? Sicher ist, dass dieses
Ma\ss{} f\"ur die Software mit zunehmender Dauer, zunehmenden Benutzern und zunehmenden
Anforderungen an Belastbarkeit verliert. Es kann nicht anders sein, denn die Unsicherheit
steigt.

Mit steigender Unsicherheit stellt sich auch die Frage der Benutzerzufriedenheit, wenn Sie
dieses Softwaresystem einsetzen. Wer kann in drei Jahren noch garantieren, dass Benutzer,
die dieses Softwaresystem einsetzen, immer noch zufrieden sind?

Was kann getan werden, um gewisse Klassen zu entwickeln, f\"ur die Einordnung und
Absch\"atzung zur Sicherheit bzw. Unsicherheit von Softwarevertr\"agen? Ist der Beruf als
Softwareentwickler noch ausreichend, um diese hoch komplexen Fragestellungen zu beantworten?

Welche anderen Berufe k\"onnten es sein? Sollten es Juristen sein? Mediziner? Mathematiker?
Physiker? Das h\"angt wohl von der Anwendung der Software und der Dom\"ane ab, in der sie
eingesetzt wird.

Es muss dar\"uber hinaus daf\"ur sensibilisiert werden, das Risiko einzugehen, ein entwickeltes
Softwaresystem einzusetzen. Eine absolute Sicherheit f\"ur ein einzusetzendes Softwaresystem
kann nach allen erfolgreichen Tests vor der Inbetriebnahme f\"ur die Produktion in der Regel
nicht garantiert werden. Das muss allen Beteiligten klar sein und diese Risiken m\"ussen
eingeplant werden. Und das ist normal und nicht unnormal. \textbf{Software hei\ss{}t Risiko.}

\subsection{Ziel dieser Arbeit}

Diese Arbeit verfolgt vier Ziele:
\begin{enumerate}
    \item Ein formales Modell der Vertragsbelastbarkeit $B(t, u, r)$ zu entwickeln und
          dessen Eigenschaften zu beweisen.
    \item Eine Klassifikation von Softwarevertr\"agen nach Unsicherheits- und
          Komplexit\"atsgrad vorzustellen.
    \item Den Satz \emph{„Software hei\ss{}t Risiko"} formal zu begr\"unden und das
          inhärente Restrisiko nach Testabschluss zu quantifizieren.
    \item Die notwendigen Qualifikationsprofile f\"ur die Bewertung von Softwarevertr\"agen
          multidisziplin\"ar zu modellieren.
\end{enumerate}

% ===========================================================================
\section{Grundlagen und Definitionen}
% ===========================================================================

\subsection{Mathematische vs.\ softwaretechnische Systeme}

\begin{definition}[Mathematisches System]
    Ein \emph{mathematisches System} $\mathcal{M} = (\mathcal{A}, \mathcal{R})$ besteht aus
    einer Menge von Axiomen $\mathcal{A}$ und einer Menge von Ableitungsregeln $\mathcal{R}$.
    F\"ur jede wohlgeformte Aussage $\varphi$ \"uber $\mathcal{M}$ gilt entweder
    $\mathcal{M} \vdash \varphi$ oder $\mathcal{M} \vdash \neg\varphi$ (sofern $\mathcal{M}$
    vollst\"andig ist).
\end{definition}

\begin{definition}[Softwaresystem]
    Ein \emph{Softwaresystem} $\mathcal{S} = (C, E, U, T)$ besteht aus einem Code-Korpus
    $C$, einer Ausf\"uhrungsumgebung $E$, einer Benutzermenge $U$ und einem Zeitparameter
    $T \in \mathbb{R}_{\geq 0}$. Die Zustandsmenge $\mathcal{Z}(\mathcal{S})$ ist im
    Allgemeinen nicht abz\"ahlbar.
\end{definition}

\begin{definition}[Formale Verifikation]
    Ein Softwaresystem $\mathcal{S}$ hei\ss{}t \emph{vollst\"andig formal verifizierbar}
    bez\"uglich einer Spezifikation $\Phi$, wenn ein Algorithmus $\mathcal{V}$ existiert
    mit:
    \[
        \mathcal{V}(\mathcal{S}, \Phi) \;\in\; \{\text{true}, \text{false}\}
        \quad \text{und} \quad
        \mathcal{V}(\mathcal{S}, \Phi) = \text{true} \;\Leftrightarrow\; \mathcal{S} \models \Phi.
    \]
\end{definition}

\begin{satz}[Unvollst\"andigkeit der Softwareverifikation]
\label{satz:unvollstaendigkeit}
    F\"ur jedes nicht-triviale Softwaresystem $\mathcal{S}$ existiert eine Spezifikation
    $\Phi$, f\"ur die $\mathcal{S}$ nicht vollst\"andig formal verifizierbar ist.
\end{satz}

\begin{beweis}
    Der Beweis folgt aus dem Halteproblem \cite{turing1936}. Sei $\mathcal{S}$ ein Turing-vollst\"andiges
    Softwaresystem. Die Spezifikation $\Phi_{\mathrm{halt}}$: \emph{„Das System terminiert
    auf allen Eingaben"} ist \"aquivalent zum Halteproblem f\"ur Turingmaschinen, das von
    Alan Turing 1936 als unentscheidbar nachgewiesen wurde \cite{turing1936}. Da jedes
    Turing-vollst\"andige System eine Turingmaschine simulieren kann, folgt: Kein Algorithmus
    $\mathcal{V}$ kann f\"ur beliebige $\mathcal{S}$ und $\Phi_{\mathrm{halt}}$ korrekt
    entscheiden. Damit ist die vollst\"andige formale Verifizierbarkeit f\"ur diese Klasse
    von Spezifikationen unm\"oglich. $\qedhere$
\end{beweis}

\subsection{Das Unsicherheitsmodell}

\begin{definition}[Unsicherheitsvektor]
    Zu einem Softwaresystem $\mathcal{S}$ und Zeitpunkt $t$ definieren wir den
    \emph{Unsicherheitsvektor}
    \[
        \mathbf{u}(t) = \bigl(u_{\mathrm{tech}}(t),\; u_{\mathrm{req}}(t),\;
                              u_{\mathrm{env}}(t),\; u_{\mathrm{usr}}(t)\bigr) \in [0,1]^4,
    \]
    wobei die Komponenten folgende Unsicherheiten messen:
    \begin{itemize}
        \item $u_{\mathrm{tech}}(t)$: technische Implementierungsunsicherheit,
        \item $u_{\mathrm{req}}(t)$: Anforderungsunsicherheit (unbekannte oder volatile
              Anforderungen),
        \item $u_{\mathrm{env}}(t)$: Umgebungsunsicherheit (physikalische Randbedingungen),
        \item $u_{\mathrm{usr}}(t)$: Benutzerunsicherheit (unbekanntes Benutzerverhalten).
    \end{itemize}
\end{definition}

\begin{definition}[Gesamtunsicherheit]
    Die \emph{Gesamtunsicherheit} $U(t)$ ist definiert als:
    \[
        U(t) \;=\; 1 - \prod_{k \in \{\mathrm{tech, req, env, usr}\}} \bigl(1 - u_k(t)\bigr).
    \]
\end{definition}

\begin{lemma}[Monotonie der Unsicherheit]
\label{lem:monotonie}
    Seien alle Komponenten $u_k(t)$ monoton nicht-fallend in $t$. Dann ist auch $U(t)$
    monoton nicht-fallend in $t$.
\end{lemma}

\begin{beweis}
    Sei $f(t) = \prod_{k}(1 - u_k(t))$. Da $u_k(t)$ nicht-fallend ist, ist $1 - u_k(t)$
    nicht-wachsend, und damit ist $f(t)$ als Produkt nicht-negativer, nicht-wachsender
    Faktoren selbst nicht-wachsend. Folglich ist $U(t) = 1 - f(t)$ nicht-fallend.
    $\qedhere$
\end{beweis}

Ein nat\"urliches Modell f\"ur das Wachstum jeder Komponente ist:
\[
    u_k(t) = 1 - e^{-\alpha_k t}, \quad \alpha_k > 0,
\]
was dem S\"attigungsverhalten realer Projekte entspricht (vgl.\ Abbildung~1).

% ===========================================================================
\section{Formales Modell der Vertragsbelastbarkeit}
% ===========================================================================

\begin{definition}[Vertragsbelastbarkeit]
\label{def:belastbarkeit}
    Sei $\mathcal{V} = (\mathcal{S}, \Phi, t_0, T_{\max})$ ein Softwarevertrag mit
    Spezifikation $\Phi$, Vertragsbeginn $t_0$ und maximaler Laufzeit $T_{\max}$.
    Der \emph{Belastbarkeitsindex}
    \[
        B\colon [0, T_{\max}] \times \mathbb{R}_{>0} \times \mathbb{R}_{>0} \;\to\; [0,1]
    \]
    ist definiert als:
    \[
        B(t, u, r) \;=\; e^{-\delta(u, r)\, t},
        \quad \text{mit} \quad
        \delta(u, r) = \delta_0 + \alpha \cdot u + \beta \cdot \log(1 + r),
    \]
    wobei $\delta_0, \alpha, \beta > 0$ systemspezifische Konstanten sind,
    $u = U(t)$ der Gesamtunsicherheitsgrad und $r$ die Anzahl der Benutzer.
\end{definition}

\begin{satz}[Verfall der Belastbarkeit]
\label{satz:verfall}
    F\"ur alle $t_1 < t_2$ gilt $B(t_1, u, r) > B(t_2, u, r)$, sofern $\delta(u,r) > 0$.
    Ferner gilt: Je gr\"o\ss{}er die Benutzerzahl $r$ und je h\"oher die Gesamtunsicherheit
    $u$, desto schneller f\"allt $B(t, u, r)$.
\end{satz}

\begin{beweis}
    \textbf{Teil 1 (Monotonier Verfall in $t$):}
    Da $\delta(u,r) > 0$, ist $e^{-\delta(u,r)\,t}$ streng monoton fallend in $t$.
    Damit gilt unmittelbar $B(t_1, u, r) > B(t_2, u, r)$ f\"ur $t_1 < t_2$. $\checkmark$

    \textbf{Teil 2 (Verst\"arkter Verfall durch $r$):}
    \[
        \frac{\partial B}{\partial r}
        = -\frac{\partial \delta}{\partial r} \cdot t \cdot B(t,u,r)
        = -\frac{\beta}{1+r} \cdot t \cdot B(t,u,r) \;<\; 0.
    \]
    Da $B > 0$ und $t > 0$, ist diese Ableitung streng negativ. Mit wachsendem $r$
    w\"achst $\delta(u,r)$, sodass $B$ schneller f\"allt. $\checkmark$

    \textbf{Teil 3 (Verst\"arkter Verfall durch $u$):}
    \[
        \frac{\partial B}{\partial u}
        = -\alpha \cdot t \cdot B(t,u,r) \;<\; 0.
    \]
    Analog folgt, dass h\"ohere Unsicherheit $u$ den Verfall von $B$ beschleunigt.
    $\qedhere$
\end{beweis}

\begin{korollar}[Keine langfristig g\"ultige Vertragssicherheit]
\label{kor:langfristig}
    F\"ur jedes $\varepsilon > 0$ und jede Systemkonfiguration mit $\delta(u,r) > 0$ existiert
    ein $T^* > 0$ mit:
    \[
        B(T^*, u, r) \;<\; \varepsilon.
    \]
    Das hei\ss{}t: Die Belastbarkeit jedes Softwarevertrags n\"ahert sich mit der Zeit
    beliebig nahe an null.
\end{korollar}

\begin{beweis}
    L\"ose $e^{-\delta t} < \varepsilon$ nach $t$ auf:
    $t > \frac{\ln(1/\varepsilon)}{\delta(u,r)} =: T^*$.
    Da $\delta(u,r) > 0$ und $\ln(1/\varepsilon) > 0$ f\"ur $\varepsilon \in (0,1)$,
    ist $T^*$ wohldefiniert und endlich. $\qedhere$
\end{beweis}

\subsection{Benutzerzufriedenheit}

\begin{definition}[Zufriedenheitsfunktion]
    Die \emph{Benutzerzufriedenheit} $S\colon [0,1] \to (0,1]$ ist definiert als:
    \[
        S(U) \;=\; \frac{1}{1 + k \cdot U}, \quad k > 0.
    \]
    Der Parameter $k$ beschreibt die Dom\"anensensitivit\"at: Bei sicherheitskritischen
    Systemen ist $k$ gro\ss{}, bei explorativen Anwendungen eher klein.
\end{definition}

\begin{satz}[Zufriedenheitsverfall]
    $S(U)$ ist streng monoton fallend in $U$. Es gilt $S(0) = 1$ und
    $\lim_{U \to \infty} S(U) = 0$.
\end{satz}

\begin{beweis}
    $\frac{\partial S}{\partial U} = -\frac{k}{(1+kU)^2} < 0$ f\"ur alle $k > 0$, $U \geq 0$.
    Die Grenzwerte ergeben sich direkt durch Einsetzen. $\qedhere$
\end{beweis}

% ===========================================================================
\section{Das inhärente Restrisiko: Software heißt Risiko}
% ===========================================================================

Dieser Abschnitt formalisiert einen der zentralen Bef\"unde dieser Arbeit: Selbst nach
vollst\"andigem Testdurchlauf und erfolgreicher Abnahme kann f\"ur den Produktionseinsatz
keine absolute Sicherheit garantiert werden. Das Risiko ist dem Medium Software strukturell
eingeschrieben.

\subsection{Testabdeckung und Restrisiko}

\begin{definition}[Testabdeckungsgrad]
    Sei $\mathcal{Z}(\mathcal{S})$ die Zustandsmenge eines Softwaresystems und $\mathcal{T}$
    eine Testmenge. Der \emph{Testabdeckungsgrad} ist definiert als:
    \[
        \tau(\mathcal{T}) \;=\; \frac{|\mathcal{T} \cap \mathcal{Z}(\mathcal{S})|}{|\mathcal{Z}(\mathcal{S})|}.
    \]
\end{definition}

\begin{satz}[Unvollst\"andigkeit endlicher Tests]
\label{satz:tests}
    Sei $\mathcal{S}$ ein nicht-triviales Softwaresystem mit $|\mathcal{Z}(\mathcal{S})| = \infty$.
    Dann gilt f\"ur jede endliche Testmenge $\mathcal{T}$:
    \[
        \tau(\mathcal{T}) \;=\; 0.
    \]
    Insbesondere kann kein endliches Testverfahren eine vollst\"andige \"Uberdeckung
    der Zustandsmenge erreichen.
\end{satz}

\begin{beweis}
    Da $|\mathcal{T}| < \infty$ und $|\mathcal{Z}(\mathcal{S})| = \infty$, gilt:
    \[
        \tau(\mathcal{T}) \;=\;
        \frac{|\mathcal{T} \cap \mathcal{Z}(\mathcal{S})|}{|\mathcal{Z}(\mathcal{S})|}
        \;\leq\;
        \frac{|\mathcal{T}|}{|\mathcal{Z}(\mathcal{S})|}
        \;=\; \frac{n}{\infty} \;=\; 0.
    \]
    Selbst wenn alle $n$ Tests bestanden werden, bleibt der Anteil getesteter Zust\"ande
    verschwindend gering gegen\"uber der unendlichen Zustandsmenge. $\qedhere$
\end{beweis}

\begin{bemerkung}
    Dijkstra formulierte diesen Sachverhalt treffend: \emph{„Program testing can be used
    to show the presence of bugs, but never to show their absence."} \cite{dijkstra1972}
    Satz~\ref{satz:tests} liefert die formale Begr\"undung dieser Aussage.
\end{bemerkung}

\begin{definition}[Restrisiko]
    Das \emph{Restrisiko} eines Softwaresystems $\mathcal{S}$ nach einer Testmenge
    $\mathcal{T}$ ist:
    \[
        \rho(\mathcal{S}, \mathcal{T}) \;=\; 1 - \tau(\mathcal{T}) \;\geq\; 1 - \frac{|\mathcal{T}|}{|\mathcal{Z}(\mathcal{S})|}.
    \]
    Im Fall $|\mathcal{Z}(\mathcal{S})| = \infty$ gilt stets $\rho(\mathcal{S}, \mathcal{T}) = 1$.
\end{definition}

\begin{korollar}[Permanenz des Restrisikos]
\label{kor:restrisiko}
    F\"ur jedes nicht-triviale Softwaresystem und jede endliche Testmenge ist das Restrisiko
    strukturell positiv:
    \[
        \rho(\mathcal{S}, \mathcal{T}) \;>\; 0.
    \]
    Es gibt keine endliche Testprozedur, die das Restrisiko auf null reduziert.
\end{korollar}

\begin{beweis}
    Folgt unmittelbar aus Satz~\ref{satz:tests}: Da $\tau(\mathcal{T}) < 1$ f\"ur jede
    endliche Testmenge, gilt $\rho = 1 - \tau > 0$. $\qedhere$
\end{beweis}

\subsection{Formaler Satz: Software heißt Risiko}

\begin{satz}[Software hei\ss{}t Risiko]
\label{satz:risiko}
    Sei $\mathcal{S}$ ein nicht-triviales, Turing-vollst\"andiges Softwaresystem.
    Dann gilt f\"ur jedes endliche Testverfahren $\mathcal{T}$ und jeden
    Produktionseinsatz mit realer Umgebung $E$:
    \[
        P\bigl(\exists\, z \in \mathcal{Z}(\mathcal{S}) \setminus \mathcal{T} \colon
        \mathcal{S}(z, E) \notin \Phi\bigr) \;>\; 0.
    \]
    Mit anderen Worten: Die Wahrscheinlichkeit, dass im Produktionsbetrieb ein nicht
    getesteter Zustand eine Spezifikationsverletzung ausl\"ost, ist strikt positiv.
    Absolute Produktionssicherheit ist daher grunds\"atzlich nicht garantierbar.
\end{satz}

\begin{beweis}
    Der Beweis gliedert sich in drei Schritte.

    \textbf{Schritt~1: Unendlichkeit der Zustandsmenge.}
    Da $\mathcal{S}$ Turing-vollst\"andig ist, kann es beliebig lange Berechnungen
    durchf\"uhren und beliebig gro\ss{}e Eingabemengen verarbeiten. Daher ist
    $|\mathcal{Z}(\mathcal{S})| = \infty$.

    \textbf{Schritt~2: Positives Restrisiko.}
    Nach Korollar~\ref{kor:restrisiko} gilt $\rho(\mathcal{S}, \mathcal{T}) > 0$, d.\,h.\
    die Menge der ungetesteten Zust\"ande $\mathcal{Z}(\mathcal{S}) \setminus \mathcal{T}$
    ist nicht-leer (tats\"achlich unendlich).

    \textbf{Schritt~3: Umgebungsvolatilit\"at.}
    Die reale Ausf\"uhrungsumgebung $E$ im Produktionsbetrieb unterscheidet sich stets
    von der Testumgebung $E_{\mathrm{test}}$ durch mindestens einen der folgenden Faktoren:
    andere Benutzeraktionen, andere Hardware-Zust\"ande, andere Netzwerkbedingungen oder
    andere Datenzust\"ande. Formal: $P(E \neq E_{\mathrm{test}}) > 0$.

    Da die Menge ungetesteter Zust\"ande nicht-leer ist und die Umgebung im Produktionsbetrieb
    von der Testumgebung abweicht, existiert mit positiver Wahrscheinlichkeit eine Kombination
    $(z, E)$ mit $z \notin \mathcal{T}$, die eine Verletzung der Spezifikation $\Phi$ ausl\"ost.
    $\qedhere$
\end{beweis}

\begin{bemerkung}[Normativit\"at des Risikos]
    Satz~\ref{satz:risiko} ist kein Versagen der Softwareentwicklung, sondern eine strukturelle
    Eigenschaft des Mediums. Das Restrisiko ist \emph{normal} und nicht \emph{abnormal}.
    Die gesellschaftliche und vertragliche Konsequenz lautet: Alle Beteiligten ---
    Auftraggeber, Auftragnehmer, Betreiber, Nutzer --- m\"ussen dieses Restrisiko bewusst
    einplanen, kommunizieren und durch geeignete Ma\ss{}nahmen (Monitoring, Fallback-Systeme,
    Versicherungen, Haftungsklauseln) adressieren.
\end{bemerkung}

\subsection{Risikobewusstsein als Vertragselement}

\begin{definition}[Risikoklausel]
    Eine \emph{Risikoklausel} $\mathcal{R}_C$ in einem Softwarevertrag ist eine Vereinbarung,
    die explizit das strukturelle Restrisiko $\rho > 0$ anerkennt und die Verantwortlichkeiten
    im Schadensfall regelt:
    \[
        \mathcal{R}_C = \bigl(\rho_{\mathrm{ack}},\; \lambda_{\text{Auftragnehmer}},\;
                               \lambda_{\text{Auftraggeber}},\; M_{\mathrm{Mitigation}}\bigr),
    \]
    wobei $\lambda$ die Haftungsanteile und $M_{\mathrm{Mitigation}}$ die vereinbarten
    Risikominderungsma\ss{}nahmen bezeichnen.
\end{definition}

\begin{hypothese}[Robustheit durch Risikotransparenz]
    Softwarevertr\"age, die eine explizite Risikoklausel $\mathcal{R}_C$ enthalten, weisen
    im Mittel einen h\"oheren Belastbarkeitsindex $B(t, u, r)$ auf als Vertr\"age ohne
    Risikoklausel, da Fehlerwartungen realistischer kalibriert sind und Eskalationen
    vermieden werden.
\end{hypothese}

% ===========================================================================
\section{Klassifikation von Softwarevertr\"agen}
% ===========================================================================

\begin{definition}[Vertragsklassen]
\label{def:klassen}
    Sei $\kappa \in [0,1]$ die \emph{technische Komplexit\"at} und
    $\nu \in [0,1]$ die \emph{Anforderungsvolatili\-t\"at} eines Softwareprojekts.
    Wir definieren vier Vertragsklassen:
    \begin{center}
    \begin{tabular}{c|c|c|l}
        \toprule
        \textbf{Klasse} & $\kappa$ & $\nu$ & \textbf{Charakteristik} \\
        \midrule
        I   & $\leq 0.5$ & $\leq 0.5$ & Stabil, formal beherrschbar \\
        II  & $> 0.5$    & $\leq 0.5$ & Technisch komplex, stabile Anforderungen \\
        III & $\leq 0.5$ & $> 0.5$    & Einfach, volatile Anforderungen \\
        IV  & $> 0.5$    & $> 0.5$    & Hoch-riskant, maximale Unsicherheit \\
        \bottomrule
    \end{tabular}
    \end{center}
\end{definition}

\begin{satz}[Belastbarkeitsordnung der Klassen]
    F\"ur die in Definition~\ref{def:klassen} definierten Klassen gilt die folgende
    Belastbarkeitsordnung zum Zeitpunkt $t > 0$:
    \[
        B_{\mathrm{I}}(t) \;\geq\; B_{\mathrm{II}}(t) \;\geq\; B_{\mathrm{III}}(t)
        \;\geq\; B_{\mathrm{IV}}(t).
    \]
\end{satz}

\begin{beweis}
    Nach Definition~\ref{def:belastbarkeit} gilt $B(t) = e^{-\delta(u,r)\,t}$ mit
    $\delta(u,r) = \delta_0 + \alpha u + \beta \log(1+r)$. Da $\alpha > 0$ und die
    Gesamtunsicherheit $u$ mit $\kappa$ und $\nu$ monoton w\"achst (mehr Komplexit\"at
    und mehr Volatili\-t\"at erh\"ohen alle Unsicherheitskomponenten), folgt:
    $\kappa_I \leq \kappa_{II} \leq \kappa_{III} \leq \kappa_{IV}$ und analog f\"ur $\nu$.
    Damit gilt $\delta_I \leq \delta_{II} \leq \delta_{III} \leq \delta_{IV}$, und
    da $e^{-x}$ streng monoton fallend in $x$ ist, folgt die behauptete Ordnung.
    $\qedhere$
\end{beweis}

\begin{beispiel}[Klassifikation realer Projekte]
    \begin{itemize}
        \item \textbf{Klasse I}: Taschenrechner-App, statische Firmenhomepage.
        \item \textbf{Klasse II}: Echtzeit-Betriebssystem f\"ur Industriesteuerungen
              (hohe Komplexit\"at, klar definierte Normen wie IEC~61508).
        \item \textbf{Klasse III}: Startup-MVP mit wochentlich wechselnden Features.
        \item \textbf{Klasse IV}: KI-gest\"utztes medizinisches Diagnosesystem,
              autonomes Fahrzeugsteuerungssystem.
    \end{itemize}
\end{beispiel}

\subsection{Informationsentropie der Anforderungen}

\begin{definition}[Anforderungsentropie]
    Sei $\mathcal{R} = \{r_1, \ldots, r_m\}$ die Anforderungsmenge mit
    Stabilit\"atswahrscheinlichkeiten $p_i$ (Wahrscheinlichkeit, dass Anforderung $r_i$
    unver\"andert bleibt). Die \emph{Anforderungsentropie} ist:
    \[
        H(\mathcal{R}) \;=\; -\sum_{i=1}^m p_i \log_2 p_i.
    \]
\end{definition}

\begin{lemma}[Entropie und Belastbarkeit]
    F\"ur steigende Anforderungsentropie $H(\mathcal{R})$ f\"allt die Belastbarkeit $B$:
    \[
        \frac{\partial B}{\partial H} \;<\; 0.
    \]
\end{lemma}

\begin{beweis}
    H\"ohere Entropie $H$ entspricht gr\"o\ss{}erer Unvorhersehbarkeit der Anforderungen,
    was die Komponente $u_{\mathrm{req}}$ erh\"oht. Nach Satz~\ref{satz:verfall} f\"allt
    $B$ mit wachsendem $u$, also gilt $\partial B / \partial H < 0$. $\qedhere$
\end{beweis}

% ===========================================================================
\section{Multidisziplin\"are Qualifikationsanforderungen}
% ===========================================================================

\begin{definition}[Kompetenzvektor]
    Das \emph{Kompetenzprofil} eines Vertragsbewerters ist ein Vektor:
    \[
        \mathbf{c} = (c_{\mathrm{SE}},\, c_{\mathrm{Jura}},\, c_{\mathrm{Med}},\,
                     c_{\mathrm{Math}},\, c_{\mathrm{Phys}},\, c_{\mathrm{Dom}}) \in [0,1]^6,
    \]
    wobei die Komponenten jeweils den normierten Kompetenzgrad in Softwareentwicklering,
    Rechtswissenschaft, Medizin/Sicherheit, Mathematik, Physik/Ingenieurwesen und
    Dom\"anenwissen messen.
\end{definition}

\begin{definition}[Mindesteignungsmatrix]
    Abh\"angig von der Vertragsklasse definieren wir eine \emph{Mindesteignungsmatrix}
    $M \in [0,1]^{4 \times 6}$:
    \[
        M \;=\;
        \begin{pmatrix}
            0.7 & 0.3 & 0.1 & 0.4 & 0.3 & 0.4 \\
            0.9 & 0.4 & 0.2 & 0.6 & 0.5 & 0.5 \\
            0.6 & 0.6 & 0.2 & 0.5 & 0.3 & 0.7 \\
            0.9 & 0.7 & 0.5 & 0.8 & 0.7 & 0.8
        \end{pmatrix},
    \]
    wobei Zeile $i$ den Mindestkompetenzvektor f\"ur Klasse $i$ darstellt.
\end{definition}

\begin{satz}[Unzul\"anglichkeit des Solo-Softwareentwickler]
\label{satz:solo}
    F\"ur Klassen~III und~IV existiert kein realistisches Einzelkompetenzprofil
    $\mathbf{c} \in [0,1]^6$ eines typischen Softwareentwickler, das alle
    Mindestanforderungen der Mindesteignungsmatrix erf\"ullt.
\end{satz}

\begin{beweis}
    Das typische Kompetenzprofil eines Softwareentwickler ist empirisch durch hohe Werte
    in $c_{\mathrm{SE}}$ und $c_{\mathrm{Math}}$ und geringe Werte in $c_{\mathrm{Jura}}$
    und $c_{\mathrm{Med}}$ gekennzeichnet \cite{ieee2014}. F\"ur Klasse~IV fordert $M$
    u.\,a.\ $c_{\mathrm{Jura}} \geq 0.7$ und $c_{\mathrm{Med}} \geq 0.5$.
    Da ein Softwareentwickler typischerweise $c_{\mathrm{Jura}} \approx 0.2$ und
    $c_{\mathrm{Med}} \approx 0.1$ aufweist, ist die Mindestanforderung verletzt.
    Folglich kann das Einzelprofil die Matrix f\"ur Klasse~IV nicht erf\"ullen.
    $\qedhere$
\end{beweis}

\begin{korollar}[Notwendigkeit interdisziplin\"arer Teams]
    F\"ur Softwareprojekte der Klassen~III und~IV ist ein interdisziplin\"ares Team
    notwendig, dessen kollektiver Kompetenzvektor $\mathbf{c}_{\mathrm{Team}}$ alle
    Eintr\"age der entsprechenden Zeile von $M$ \"uberschreitet.
\end{korollar}

% ===========================================================================
\section{Visualisierung der Ergebnisse}
% ===========================================================================

Alle folgenden Abbildungen wurden mit Python/Matplotlib erzeugt und liegen als separate
PDF-Dateien vor.

\subsection{Abbildung 1: Wachstum der Unsicherheit}

\begin{figure}[h!]
    \centering
    \includegraphics[width=0.85\textwidth]{plots/plot1_unsicherheit.pdf}
    \caption{Wachstum der technischen Unsicherheit $U_{\mathrm{tech}}(t)$,
    der Anforderungsunsicherheit $U_{\mathrm{req}}(t)$ und der Gesamtunsicherheit
    $U_{\mathrm{ges}}(t)$ \"uber die normierte Projektlaufzeit $t$. Alle Kurven folgen
    einem S\"attigungsmodell $U_k(t) = 1 - e^{-\alpha_k t}$ mit domänenspezifischen
    Zerfallsraten $\alpha_k$. Die Gesamtunsicherheit \"uberschreitet bereits bei
    $t \approx 7$ den Wert $0.9$, was einer praktisch nicht mehr kontrollierbaren
    Vertragssituation entspricht. Dies visualisiert formal die Aussage von
    Satz~\ref{satz:verfall} und Lemma~\ref{lem:monotonie}.}
    \label{fig:unsicherheit}
\end{figure}

\subsection{Abbildung 2: Zeitlicher Verfall des Belastbarkeitsindex}

\begin{figure}[h!]
    \centering
    \includegraphics[width=0.85\textwidth]{plots/plot2_belastbarkeit.pdf}
    \caption{Zeitlicher Verlauf des Belastbarkeitsindex $B(t)$ gem\"a\ss{}
    Definition~\ref{def:belastbarkeit} f\"ur drei Szenarien: idealer Verlauf
    $B_{\mathrm{ideal}}(t) = e^{-\delta t}$, realistischer Verlauf mit st\"ochastischen
    Schwankungen (z.\,B.\ durch Anforderungs\"anderungen), und Minimalverlauf bei
    erh\"ohter Benutzerzahl. Das Unsicherheitsband zwischen idealem und minimalem
    Verlauf vergr\"o\ss{}ert sich mit der Zeit. Alle Kurven konvergieren gegen null,
    entsprechend Korollar~\ref{kor:langfristig}.}
    \label{fig:belastbarkeit}
\end{figure}

\subsection{Abbildung 3: Benutzerzufriedenheit als Funktion der Unsicherheit}

\begin{figure}[h!]
    \centering
    \includegraphics[width=0.95\textwidth]{plots/plot3_zufriedenheit.pdf}
    \caption{Benutzerzufriedenheit $S(U)$ f\"ur verschiedene Dom\"anensensitivit\"aten
    $k$. Linke Teilabbildung ($k=1$): Standardanwendung. Rechte Teilabbildung ($k=3$):
    sicherheitskritische Anwendung. Der Vergleich zeigt, dass die Zufriedenheit in
    sicherheitskritischen Dom\"anen bereits bei moderater Unsicherheit $U \approx 0.3$
    auf unter $50\,\%$ f\"allt. Die quadratische Variante modelliert nichtlineares
    Toleranzverhalten der Benutzer.}
    \label{fig:zufriedenheit}
\end{figure}

\subsection{Abbildung 4: Klassifikationsmatrix}

\begin{figure}[h!]
    \centering
    \includegraphics[width=0.65\textwidth]{plots/plot4_klassifikation.pdf}
    \caption{Klassifikationsmatrix f\"ur Softwarevertr\"age nach technischer Komplexit\"at
    $\kappa$ und Anforderungsvolatili\-t\"at $\nu$. Die vier Klassen aus
    Definition~\ref{def:klassen} sind farblich kodiert: Klasse~I (gr\"un, niedrigstes
    Risiko) bis Klasse~IV (rot, h\"ochstes Risiko). Die Farbkodierung entspricht
    intuitiv einem Ampelsystem, das Projektbeteiligte bei der Risikoeinsch\"atzung
    unterst\"utzt.}
    \label{fig:klassifikation}
\end{figure}

\subsection{Abbildung 5: Anforderungsentropie}

\begin{figure}[h!]
    \centering
    \includegraphics[width=0.85\textwidth]{plots/plot5_entropie.pdf}
    \caption{Zeitlicher Verlauf der normierten Anforderungsentropie $H(t)$ mit
    unterer Schranke. Die Entropie zeigt einen Gipfel in der Projektmitte --- typisch
    f\"ur Projekte, in denen Anforderungen zun\"achst wachsen (Discovery-Phase) und
    sp\"ater durch zunehmende Kl\"arung wieder fallen. Das Unsicherheitsband zeigt,
    dass selbst die untere Schranke nie auf null f\"allt, was das permanente Restrisiko
    aus Korollar~\ref{kor:restrisiko} visualisiert.}
    \label{fig:entropie}
\end{figure}

\subsection{Abbildung 6: Multidisziplin\"ares Kompetenzprofil}

\begin{figure}[h!]
    \centering
    \includegraphics[width=0.72\textwidth]{plots/plot6_kompetenz.pdf}
    \caption{Radardiagramm der Kompetenzprofile f\"ur vier Berufsgruppen:
    Softwareentwickler, Jurist, Mathematiker und ein interdisziplin\"ares Team.
    Deutlich erkennbar ist, dass kein Einzelprofil alle sechs Kompetenzdimensionen
    gleichzeitig abdeckt --- die formale Grundlage f\"ur Satz~\ref{satz:solo}.
    Nur das interdisziplin\"are Team n\"ahert sich einem gleichm\"a\ss{}ig hohen
    Kompetenzprofil.}
    \label{fig:kompetenz}
\end{figure}

% ===========================================================================
\section{Konsequenzen f\"ur die UML-Modellierung und das UML-Risikoprofil}
\label{sec:uml}
% ===========================================================================

Die bisherigen Ergebnisse haben gezeigt, dass Softwarevertr\"age inhärent unsicher sind
und dass das Restrisiko $\rho > 0$ strukturell nicht eliminierbar ist. Diese Erkenntnis
hat unmittelbare Konsequenzen f\"ur die \emph{Modellierung} von Softwaresystemen: Wenn
Risiko dem System strukturell innewohnt, m\"ussen auch die Modelle, mit denen das System
beschrieben und entworfen wird, in der Lage sein, dieses Risiko zu repr\"asentieren und
zu quantifizieren. Die \emph{Unified Modeling Language} (UML) in ihrer Standardform
\cite{omg2017} bietet hierf\"ur keine eingebauten Mechanismen. In diesem Kapitel leiten
wir formal ab, welche Konsequenzen die Ergebnisse der vorangegangenen Abschnitte f\"ur
die UML-Modellierung haben, und definieren ein \emph{UML-Risikoprofil} \textsf{RiskProfile},
das die quantitative Annotation von Risikoma\ss{}zahlen direkt in UML-Diagrammen erm\"oglicht.

% ---------------------------------------------------------------------------
\subsection{Grundlagen: UML-Profile als Erweiterungsmechanismus}
% ---------------------------------------------------------------------------

\begin{definition}[UML-Profil]
\label{def:uml-profil}
    Ein \emph{UML-Profil} $\Pi = (\mathcal{ST}, \mathcal{TV}, \mathcal{C})$ besteht
    aus einer Menge von Stereotypen $\mathcal{ST}$, einer Menge von \emph{Tagged Values}
    $\mathcal{TV}$ und einer Menge von Einschr\"ankungen $\mathcal{C}$
    (OCL-Constraints\footnote{OCL: Object Constraint Language, standardisierter Teil
    von UML~2.5 \cite{omg2017}.}). Ein Stereotyp $s \in \mathcal{ST}$ erweitert genau
    ein UML-Metaklassenelement $m \in \mathcal{M}_{\mathrm{UML}}$, notiert als
    $s \triangleright m$.
\end{definition}

\begin{definition}[UML-Metamodell]
    Das \emph{UML-Metamodell} $\mathcal{M}_{\mathrm{UML}}$ ist die Menge aller
    abstrakten Syntaxelemente der UML~2.5, darunter insbesondere:
    \begin{align*}
        \mathcal{M}_{\mathrm{UML}} \;\supseteq\;
        \{\,&\texttt{Class},\; \texttt{Component},\; \texttt{Interface},\;
               \texttt{Operation},\\
            &\texttt{Message},\; \texttt{Lifeline},\;
               \texttt{Dependency},\; \texttt{Association}\,\}.
    \end{align*}
\end{definition}

UML-Profile sind der von der OMG standardisierte leichtgewichtige Erweiterungsmechanismus
\cite{omg2017}. Im Unterschied zur Neudefinition von Metamodell-Elementen bleibt die
Kompatibilit\"at zur Basis-UML erhalten. Dies ist f\"ur das \textsf{RiskProfile}
entscheidend: Es soll mit bestehenden UML-Werkzeugen kompatibel sein.

% ---------------------------------------------------------------------------
\subsection{Definition des UML-Risikoprofils \textsf{RiskProfile}}
% ---------------------------------------------------------------------------

\begin{definition}[Risikoma\ss{}zahl]
\label{def:risikozahl}
    Die \emph{Risikoma\ss{}zahl} eines UML-Modellelements $e \in \mathcal{M}_{\mathrm{UML}}$
    ist ein Tupel:
    \[
        \omega(e) \;=\; \bigl(\rho_e,\; u_e,\; \lambda_e,\; \kappa_e\bigr)
        \;\in\; [0,1]^4,
    \]
    wobei
    \begin{itemize}
        \item $\rho_e \in [0,1]$: das dem Element zugeordnete Restrisiko
              (gem\"a\ss{} Definition aus Abschnitt~4),
        \item $u_e \in [0,1]$: der lokale Unsicherheitsgrad des Elements,
        \item $\lambda_e \in [0,1]$: die Laufzeitkritikalit\"at
              (wie stark beeinflusst ein Fehler in $e$ den Gesamtbetrieb),
        \item $\kappa_e \in [0,1]$: die Kopplungskomplexit\"at des Elements
              (normierte strukturelle Kopplung, z.\,B.\ Fan-in/Fan-out-Verh\"altnis).
    \end{itemize}
\end{definition}

\begin{definition}[UML-Risikoprofil \textsf{RiskProfile}]
\label{def:riskprofile}
    Das \emph{UML-Risikoprofil} $\Pi_R = (\mathcal{ST}_R, \mathcal{TV}_R, \mathcal{C}_R)$
    ist definiert durch folgende Stereotypen und Tagged Values:

    \medskip
    {\small
    \begin{center}
    \renewcommand{\arraystretch}{1.25}
    \begin{tabular}{p{3.8cm}p{2.4cm}p{6.5cm}}
        \toprule
        \textbf{Stereotyp} & \textbf{Erweitert} & \textbf{Tagged Values} \\
        \midrule
        \texttt{\guillemotleft riskElement\guillemotright}
            & \texttt{Element}
            & \texttt{residualRisk: Real},
              \texttt{uncertainty: Real},
              \texttt{criticality: Real},
              \texttt{coupling: Real} \\[2pt]
        \texttt{\guillemotleft riskClass\guillemotright}
            & \texttt{Class}
            & erbt von \texttt{riskElement};
              \texttt{stateSpaceEst: String} \\[2pt]
        \texttt{\guillemotleft riskComponent\guillemotright}
            & \texttt{Component}
            & erbt von \texttt{riskElement};
              \texttt{deployRisk: Real} \\[2pt]
        \texttt{\guillemotleft riskMessage\guillemotright}
            & \texttt{Message}
            & erbt von \texttt{riskElement};
              \texttt{failProb: Real} \\[2pt]
        \texttt{\guillemotleft riskLifeline\guillemotright}
            & \texttt{Lifeline}
            & erbt von \texttt{riskElement};
              \texttt{availabilityTarget: Real} \\[2pt]
        \texttt{\guillemotleft riskDependency\guillemotright}
            & \texttt{Dependency}
            & erbt von \texttt{riskElement};
              \texttt{propagationFactor: Real} \\
        \bottomrule
    \end{tabular}
    \end{center}
    }

    \medskip\noindent
    Die Einschr\"ankungsmenge $\mathcal{C}_R$ enth\"alt die OCL-Constraints:
    \begin{align*}
        C_1&: \quad \forall\, e \;\text{mit Stereotyp}\;
               \texttt{riskElement}\colon
               \rho_e, u_e, \lambda_e, \kappa_e \in [0,1], \\
        C_2&: \quad \texttt{failProb} \;\leq\; \rho_e
               \quad (\text{Aufruf-Fehlerwahrscheinlichkeit} \leq \text{Restrisiko}), \\
        C_3&: \quad \texttt{propagationFactor} \;\in\; [0,1], \\
        C_4&: \quad \texttt{availabilityTarget} \;\in\; (0,1].
    \end{align*}
\end{definition}

% ---------------------------------------------------------------------------
\subsection{Formale Verkn\"upfung von \textsf{RiskProfile} und\\Belastbarkeitsmodell}
% ---------------------------------------------------------------------------

\begin{satz}[Konsistenz von \textsf{RiskProfile} und Belastbarkeitsindex]
\label{satz:uml-konsistenz}
    Sei $\mathcal{S}$ ein Softwaresystem, dessen UML-Modell mit dem Profil
    $\Pi_R$ annotiert ist. Sei $E = \{e_1, \ldots, e_n\}$ die Menge aller
    annotierten Modellelemente. Dann gilt:
    \[
        B(t, u, r) \;\leq\; 1 - \max_{i=1,\ldots,n} \lambda_{e_i} \cdot \rho_{e_i},
    \]
    d.\,h.\ der Belastbarkeitsindex des Gesamtsystems ist nach oben beschr\"ankt
    durch das risikokritischste annotierte Element.
\end{satz}

\begin{beweis}
    Sei $e^* = \arg\max_i \lambda_{e_i} \cdot \rho_{e_i}$ das kritischste Element.
    Nach Korollar~\ref{kor:restrisiko} gilt $\rho_{e^*} > 0$. Da $\lambda_{e^*}$
    die Laufzeitkritikalit\"at misst (Anteil des Gesamtsystemversagens, der durch
    einen Fehler in $e^*$ erkl\"art wird), gilt:
    \[
        P(\text{Systemausfall}) \;\geq\; \lambda_{e^*} \cdot \rho_{e^*} \;=:\; p^*.
    \]
    Der Belastbarkeitsindex ist definiert als $B = e^{-\delta t}$, wobei
    $\delta \geq -\frac{1}{t}\ln(1 - p^*)$ gelten muss, damit der Ausfall in
    $[0,t]$ mit Wahrscheinlichkeit $p^*$ eintritt. Daher:
    \[
        B(t,u,r) \;\leq\; 1 - p^* \;=\; 1 - \max_i \lambda_{e_i} \cdot \rho_{e_i}.
        \quad\qedhere
    \]
\end{beweis}

\begin{lemma}[Monotonie der Risikoannotation]
\label{satz:uml-monoton}
    Seien $e, e' \in E$ zwei Modellelemente mit $\kappa_e \leq \kappa_{e'}$ und
    $u_e \leq u_{e'}$. Dann gilt:
    \[
        \omega(e) \;\preceq\; \omega(e'),
    \]
    wobei $\preceq$ die komponentenweise Halbordnung auf $[0,1]^4$ bezeichnet.
    Ferner impliziert $\omega(e) \preceq \omega(e')$, dass Element $e'$ eine
    h\"ohere Priorit\"at bei Risikobehandlungsma\ss{}nahmen erh\"alt.
\end{lemma}

\begin{beweis}
    Da $\kappa$ die Kopplungskomplexit\"at und $u$ den Unsicherheitsgrad misst,
    und beide in das Restrisiko $\rho$ eingehen gem\"a\ss{}
    $\rho_e = 1 - (1-u_e)(1-\kappa_e \cdot c)$ f\"ur eine Konstante $c > 0$,
    folgt: $u_e \leq u_{e'}$ und $\kappa_e \leq \kappa_{e'}$ impliziert
    $\rho_e \leq \rho_{e'}$. Damit gilt die komponentenweise Ungleichung
    $\omega(e) \preceq \omega(e')$. Die Priorit\"atsordnung folgt direkt aus
    der Definition der Risikoma\ss{}zahl. $\qedhere$
\end{beweis}

% ---------------------------------------------------------------------------
\subsection{Klassendiagramme: Strukturelles Risiko}
% ---------------------------------------------------------------------------

Klassendiagramme beschreiben die statische Struktur eines Softwaresystems. Die
Risikorelevanz ergibt sich hier prim\"ar aus der Kopplung zwischen Klassen und
der Gr\"o\ss{}e des Zustandsraums einzelner Klassen.

\begin{definition}[Strukturelles Risiko einer Klasse]
\label{def:strukturrisiko}
    Sei $C$ eine UML-Klasse mit $n_A$ Attributen, $n_O$ Operationen und
    $\deg(C)$ dem Grad im Assoziationsgraphen (Fan-in + Fan-out). Das
    \emph{strukturelle Risiko} von $C$ ist:
    \[
        \rho^{\mathrm{struct}}(C) \;=\; 1 - e^{-\gamma \cdot \deg(C)} \cdot
        \frac{1}{1 + \mu \cdot n_A \cdot n_O},
    \]
    mit Skalierungskonstanten $\gamma, \mu > 0$.
\end{definition}

\begin{satz}[Monotonie des strukturellen Risikos]
\label{satz:strukturrisiko-monoton}
    $\rho^{\mathrm{struct}}(C)$ ist streng monoton wachsend in $\deg(C)$, $n_A$ und
    $n_O$.
\end{satz}

\begin{beweis}
    \textbf{Monotonie in $\deg(C)$:} Da $e^{-\gamma \cdot \deg(C)}$ streng monoton
    fallend in $\deg(C)$ ist, ist $1 - e^{-\gamma \cdot \deg(C)}$ streng monoton
    wachsend. Der Faktor $\frac{1}{1+\mu n_A n_O} \leq 1$ skaliert diesen Beitrag,
    \"andert aber nicht die Monotoniezusage, da er nicht von $\deg(C)$ abh\"angt.

    \textbf{Monotonie in $n_A$ und $n_O$:} Der Term $\frac{1}{1 + \mu n_A n_O}$
    ist streng monoton fallend in $n_A$ und $n_O$ (da der Nenner w\"achst). Damit
    ist $\rho^{\mathrm{struct}}(C) = 1 - e^{-\gamma \deg(C)} \cdot \frac{1}{1+\mu n_A n_O}$
    streng wachsend in $n_A$ und $n_O$, denn $e^{-\gamma \deg(C)} > 0$ und ein
    Subtrahend, der f\"allt, ergibt einen Ausdruck, der steigt. $\qedhere$
\end{beweis}

\begin{korollar}[Risikopriorisierung im Klassendiagramm]
\label{kor:klasse-prio}
    Im mit $\Pi_R$ annotierten Klassendiagramm erhalten Klassen mit hohem Kopplungsgrad
    $\deg(C)$ und vielen Attribut-Operationen-Kombinationen die h\"ochsten
    \texttt{residualRisk}-Werte und sind vorrangig durch Entkopplungsma\ss{}nahmen
    (z.\,B.\ Einf\"uhrung von Fassaden oder Interfaces) zu behandeln.
\end{korollar}

\begin{beispiel}[Risikoannotation einer Klasse]
    Betrachte eine Klasse \texttt{OrderManager} mit $n_A = 12$, $n_O = 18$,
    $\deg(\texttt{OrderManager}) = 9$ (hohe Kopplung). Mit $\gamma = 0.15$ und
    $\mu = 0.002$ ergibt sich:
    \[
        \rho^{\mathrm{struct}}(\texttt{OrderManager})
        = 1 - e^{-0.15 \cdot 9} \cdot \frac{1}{1 + 0.002 \cdot 12 \cdot 18}
        \approx 1 - 0.259 \cdot 0.698 \approx 0.819.
    \]
    Im UML-Diagramm w\"are diese Klasse mit
    \texttt{\guillemotleft{}riskClass\guillemotright{} residualRisk=0.82, coupling=0.9}
    annotiert und w\"are damit ein vorrangiger Refaktorierungskandidat.
\end{beispiel}

% ---------------------------------------------------------------------------
\subsection{Sequenzdiagramme: Dynamisches Interaktionsrisiko}
% ---------------------------------------------------------------------------

Sequenzdiagramme beschreiben das zeitliche Zusammenspiel von Objekten und Nachrichten.
Das Risiko manifestiert sich hier in der Fehlerwahrscheinlichkeit einzelner
Nachrichtenaufrufe und deren Fortpflanzung entlang der Aufrufkette.

\begin{definition}[Interaktionsrisiko einer Nachrichtenkette]
\label{def:interaktionsrisiko}
    Sei $\sigma = (m_1, m_2, \ldots, m_k)$ eine Sequenz von Nachrichten in einem
    UML-Sequenzdiagramm. Das \emph{Interaktionsrisiko} von $\sigma$ ist:
    \[
        \rho^{\mathrm{seq}}(\sigma) \;=\; 1 - \prod_{i=1}^{k}
        \bigl(1 - \rho_{m_i}\bigr),
    \]
    wobei $\rho_{m_i}$ das \texttt{failProb}-Tagged-Value der Nachricht $m_i$ im
    \textsf{RiskProfile} bezeichnet.
\end{definition}

\begin{satz}[Superadditivit\"at des Interaktionsrisikos]
\label{satz:seq-superadditiv}
    F\"ur alle Nachrichtensequenzen $\sigma = (m_1, \ldots, m_k)$ mit
    $\rho_{m_i} > 0$ gilt:
    \[
        \rho^{\mathrm{seq}}(\sigma) \;>\; \max_{i=1,\ldots,k} \rho_{m_i}.
    \]
    Das Gesamtrisiko einer Nachrichtenkette \"ubersteigt stets das Maximum der
    Einzelrisiken.
\end{satz}

\begin{beweis}
    Sei $\rho^* = \max_i \rho_{m_i}$. Dann gilt:
    \[
        \prod_{i=1}^k (1 - \rho_{m_i}) \;\leq\; (1 - \rho^*) \cdot
        \prod_{i \neq i^*} (1 - \rho_{m_i}) \;<\; (1 - \rho^*),
    \]
    da jeder Faktor $(1 - \rho_{m_i}) < 1$ f\"ur $\rho_{m_i} > 0$. Daher:
    \[
        \rho^{\mathrm{seq}}(\sigma) = 1 - \prod_{i=1}^k(1-\rho_{m_i})
        > 1 - (1-\rho^*) = \rho^*.
        \quad\qedhere
    \]
\end{beweis}

\begin{korollar}[Risikokritischer Pfad im Sequenzdiagramm]
\label{kor:kritpfad}
    Sei $\Sigma$ die Menge aller m\"oglichen Nachrichtensequenzen in einem
    Sequenzdiagramm. Der \emph{risikokritische Pfad} ist:
    \[
        \sigma^* \;=\; \arg\max_{\sigma \in \Sigma}\; \rho^{\mathrm{seq}}(\sigma).
    \]
    Dieser Pfad ist vorrangig durch Fehlerbehandlungsmechanismen (z.\,B.\ Try-Catch,
    Retry-Logik, Circuit-Breaker) abzusichern.
\end{korollar}

\begin{beispiel}[Interaktionsrisiko einer Login-Sequenz]
    Betrachte eine Login-Sequenz mit f\"unf Nachrichten:
    \[
        \rho_{m_1} = 0.02,\quad \rho_{m_2} = 0.05,\quad \rho_{m_3} = 0.03,
        \quad \rho_{m_4} = 0.08,\quad \rho_{m_5} = 0.01.
    \]
    Das Interaktionsrisiko betr\"agt:
    \[
        \rho^{\mathrm{seq}}(\sigma)
        = 1 - (0.98 \cdot 0.95 \cdot 0.97 \cdot 0.92 \cdot 0.99)
        = 1 - 0.819 \approx 0.181.
    \]
    Obwohl kein Einzelrisiko $0.1$ \"ubersteigt, liegt das Gesamtrisiko der
    Sequenz bei $18{,}1\,\%$ --- ein Beleg f\"ur Satz~\ref{satz:seq-superadditiv}.
\end{beispiel}

% ---------------------------------------------------------------------------
\subsection{Komponentendiagramme: Einsatz- und Schnittstellenrisiko}
% ---------------------------------------------------------------------------

Komponentendiagramme beschreiben die Zerlegung eines Softwaresystems in
deploybare Einheiten und deren Schnittstellen. Das Risiko manifestiert sich
hier in der Abh\"angigkeitspropagation zwischen Komponenten.

\begin{definition}[Komponentenrisikograph]
\label{def:kompgraph}
    Sei $\mathcal{G} = (V, E_K)$ ein gerichteter Graph, wobei $V$ die Menge der
    UML-Komponenten und $E_K \subseteq V \times V$ die Abh\"angigkeitsbeziehungen
    (Use-Relationships) bezeichnen. Jeder Kante $(c_i, c_j) \in E_K$ ist ein
    \emph{Propagationsfaktor} $\pi_{ij} \in [0,1]$ zugeordnet
    (\texttt{propagationFactor} im \textsf{RiskProfile}).
\end{definition}

\begin{satz}[Risikopropagation in Komponentengraphen]
\label{satz:komp-propagation}
    Sei $c_0$ eine Wurzelkomponente mit Restrisiko $\rho_{c_0}$ und
    $(c_0, c_1, \ldots, c_\ell)$ ein gerichteter Pfad der L\"ange $\ell$ in
    $\mathcal{G}$. Das propagierte Risiko an der Endkomponente $c_\ell$ ist:
    \[
        \rho^{\mathrm{prop}}(c_\ell) \;=\; \rho_{c_\ell} +
        \left(1 - \rho_{c_\ell}\right) \cdot \rho_{c_0} \cdot
        \prod_{k=0}^{\ell-1} \pi_{c_k c_{k+1}}.
    \]
\end{satz}

\begin{beweis}
    Das Ereignis „Ausfall von $c_\ell$" tritt ein entweder durch:
    \begin{enumerate}[label=(\roman*)]
        \item Eigenversagen von $c_\ell$ mit Wahrscheinlichkeit $\rho_{c_\ell}$, oder
        \item Versagen durch Propagation aus $c_0$: Wahrscheinlichkeit
              $\rho_{c_0} \cdot \prod_{k=0}^{\ell-1} \pi_{c_k c_{k+1}}$.
    \end{enumerate}
    Da beide Ereignisse nicht notwendig disjunkt sind, gilt mit der
    Inklusion-Exklusion-Formel:
    \begin{align*}
        \rho^{\mathrm{prop}}(c_\ell)
        &= \rho_{c_\ell} + \rho_{c_0} \prod_k \pi_{c_k c_{k+1}}
           - \rho_{c_\ell} \cdot \rho_{c_0} \prod_k \pi_{c_k c_{k+1}} \\
        &= \rho_{c_\ell} + \bigl(1 - \rho_{c_\ell}\bigr) \cdot
           \rho_{c_0} \prod_k \pi_{c_k c_{k+1}}.
        \quad\qedhere
    \end{align*}
\end{beweis}

\begin{korollar}[Maximalrisiko durch Propagation]
\label{kor:komp-max}
    F\"ur $\pi_{ij} = 1$ (vollst\"andige Propagation) und $\rho_{c_0} = 1$ gilt:
    \[
        \rho^{\mathrm{prop}}(c_\ell) \;=\; 1.
    \]
    Vollst\"andige Abh\"angigkeit bei maximalem Quellrisiko f\"uhrt zu absolutem
    Ausfallrisiko der Zielkomponente --- eine formale Begr\"undung f\"ur das Prinzip
    der losen Kopplung (\emph{loose coupling}) in der Softwarearchitektur.
\end{korollar}

\begin{beweis}
    Einsetzen in die Formel aus Satz~\ref{satz:komp-propagation}:
    $\rho^{\mathrm{prop}}(c_\ell) = \rho_{c_\ell} + (1-\rho_{c_\ell}) \cdot 1 \cdot 1 = 1$.
    $\qedhere$
\end{beweis}

% ---------------------------------------------------------------------------
\subsection{Weitere relevante Diagrammtypen}
% ---------------------------------------------------------------------------

\subsubsection{Zustandsdiagramme}

Zustandsdiagramme (\emph{State Machine Diagrams}) modellieren das Verhalten von
Objekten als endliche Automaten. Die Risikorelevanz ergibt sich aus der Anzahl
erreichbarer Zust\"ande und den Transitionsbedingungen.

\begin{definition}[Zustandsraumrisiko]
\label{def:zustandsrisiko}
    Sei $\mathcal{A} = (Q, \Sigma, \delta, q_0, F)$ ein endlicher Automat, der
    das Verhalten einer Klasse $C$ modelliert. Das \emph{Zustandsraumrisiko} ist:
    \[
        \rho^{\mathrm{state}}(\mathcal{A}) \;=\; 1 - \frac{|Q_{\mathrm{getestet}}|}{|Q|},
    \]
    wobei $Q_{\mathrm{getestet}} \subseteq Q$ die Menge der durch Testf\"alle
    abgedeckten Zust\"ande bezeichnet.
\end{definition}

\begin{satz}[Zustandsraumrisiko und Testabdeckung]
\label{satz:zustand-test}
    F\"ur jede unvollst\"andige Testmenge gilt
    $\rho^{\mathrm{state}}(\mathcal{A}) > 0$.
    Ferner w\"achst $\rho^{\mathrm{state}}$ monoton mit $|Q|$ bei konstanter
    Testmengengr\"o\ss{}e.
\end{satz}

\begin{beweis}
    Da $|Q_{\mathrm{getestet}}| \leq |Q_{\mathrm{test}}| < |Q|$ bei unvollst\"andiger
    Testmenge, gilt $\frac{|Q_{\mathrm{getestet}}|}{|Q|} < 1$, also
    $\rho^{\mathrm{state}} > 0$. Sei $|Q_{\mathrm{getestet}}|$ konstant und $|Q|$
    wachsend. Dann f\"allt der Bruch $\frac{|Q_{\mathrm{getestet}}|}{|Q|}$ streng
    monoton, und $\rho^{\mathrm{state}} = 1 - \frac{|Q_{\mathrm{getestet}}|}{|Q|}$
    w\"achst streng monoton. $\qedhere$
\end{beweis}

\subsubsection{Aktivit\"atsdiagramme}

Aktivit\"atsdiagramme modellieren Arbeitsabl\"aufe und Kontrollfl\"usse. Das Risiko
ergibt sich aus der Komplexit\"at der Verzweigungslogik.

\begin{definition}[Kontrollflussrisiko]
\label{def:kfrisiko}
    Sei $G_A = (N, E_A)$ der Kontrollfluss-Graph eines Aktivit\"atsdiagramms mit
    $|N|$ Knoten und $|E_A|$ Kanten. Das \emph{Kontrollflussrisiko} basiert auf der
    zyklomatischen Komplexit\"at \cite{mccabe1976}:
    \[
        M(G_A) \;=\; |E_A| - |N| + 2,
    \]
    und das Kontrollflussrisiko ist:
    \[
        \rho^{\mathrm{ctrl}}(G_A) \;=\; 1 - e^{-\eta \cdot M(G_A)}, \quad \eta > 0.
    \]
\end{definition}

\begin{lemma}[Kontrollflussrisiko und zyklomatische Komplexit\"at]
\label{lem:kf-zyklom}
    $\rho^{\mathrm{ctrl}}(G_A)$ ist streng monoton wachsend in $M(G_A)$.
    F\"ur $M(G_A) \leq 10$ (McCabe-Grenze f\"ur wartbare Komplexit\"at \cite{mccabe1976})
    gilt $\rho^{\mathrm{ctrl}}(G_A) < 1 - e^{-10\eta}$.
\end{lemma}

\begin{beweis}
    Da $e^{-\eta M}$ streng monoton fallend in $M > 0$ ist, folgt die erste Aussage
    unmittelbar. Die Schranke folgt durch Einsetzen von $M(G_A) = 10$:
    $\rho^{\mathrm{ctrl}} = 1 - e^{-10\eta} < 1$. $\qedhere$
\end{beweis}

\subsubsection{Anwendungsfall- und Verteilungsdiagramme}

\begin{bemerkung}[Anwendungsfalldiagramme]
    Anwendungsfalldiagramme (\emph{Use Case Diagrams}) geben prim\"ar eine
    Außenperspektive auf das System. Im \textsf{RiskProfile} werden Use Cases mit
    \texttt{\guillemotleft{}riskElement\guillemotright{}} annotiert, wobei
    \texttt{uncertainty} die Unklarheit der Anforderungsdefinition und
    \texttt{criticality} die gesch\"aftliche Kritikalit\"at repr\"asentiert. Ein
    Use Case mit hoher \texttt{uncertainty} korrespondiert direkt mit einer hohen
    Komponente $u_{\mathrm{req}}$ im Unsicherheitsvektor aus Abschnitt~2.
\end{bemerkung}

\begin{bemerkung}[Verteilungsdiagramme]
    Verteilungsdiagramme (\emph{Deployment Diagrams}) zeigen die physikalische
    Verteilung von Softwarekomponenten auf Knoten. Das \texttt{deployRisk}-Tagged-Value
    des Stereotyps \texttt{\guillemotleft{}riskComponent\guillemotright{}} modelliert
    hier das Risiko aus Umgebungsunsicherheit $u_{\mathrm{env}}$ --- also die
    Wahrscheinlichkeit, dass der Zielknoten nicht in der Konfiguration vorliegt,
    f\"ur die das System entwickelt wurde.
\end{bemerkung}

% ---------------------------------------------------------------------------
\subsection{Aggregation: Systemweites Risikomodell}
% ---------------------------------------------------------------------------

\begin{definition}[Systemweites Risiko eines UML-Modells]
\label{def:systemrisiko}
    Sei $\mathcal{D} = \{D_1, \ldots, D_p\}$ eine Sammlung von UML-Diagrammen
    (Klassen-, Sequenz-, Komponenten-, Zustands- und Aktivit\"atsdiagramme),
    alle annotiert mit $\Pi_R$. Das \emph{systemweite Risiko} ist:
    \[
        \rho^{\mathrm{sys}}(\mathcal{D}) \;=\; 1 - \prod_{j=1}^{p}
        \bigl(1 - w_j \cdot \overline{\rho}(D_j)\bigr),
    \]
    wobei $w_j \in [0,1]$ das Gewicht des Diagrammtyps (mit $\sum_j w_j = 1$) und
    \[
        \overline{\rho}(D_j) \;=\; \frac{1}{|E_j|}
        \sum_{e \in E_j} \lambda_e \cdot \rho_e
    \]
    das gewichtete mittlere Elementrisiko des Diagramms $D_j$ bezeichnet, mit
    $E_j$ als Menge der annotierten Elemente in $D_j$.
\end{definition}

\begin{satz}[Risikoaggregationstheorem]
\label{satz:uml-aggregation}
    Das systemweite Risiko $\rho^{\mathrm{sys}}(\mathcal{D})$ besitzt folgende
    Eigenschaften:
    \begin{enumerate}[label=(\roman*)]
        \item $\rho^{\mathrm{sys}}(\mathcal{D}) > 0$ f\"ur alle nicht-leeren
              Diagrammmengen mit $\overline{\rho}(D_j) > 0$.
        \item $\rho^{\mathrm{sys}}(\mathcal{D})$ ist streng monoton wachsend in
              jedem $\overline{\rho}(D_j)$.
        \item Es gilt: $\rho^{\mathrm{sys}}(\mathcal{D}) < 1$ genau dann, wenn
              $\overline{\rho}(D_j) < 1$ f\"ur alle $j$.
        \item $\rho^{\mathrm{sys}}(\mathcal{D}) \geq \max_j w_j \cdot
              \overline{\rho}(D_j)$.
    \end{enumerate}
\end{satz}

\begin{beweis}
    \textbf{(i):} F\"ur $\overline{\rho}(D_j) > 0$ gilt
    $w_j \cdot \overline{\rho}(D_j) > 0$ (da $w_j > 0$), also
    $(1 - w_j\overline{\rho}(D_j)) < 1$. Das Produkt $\prod_j (1-w_j\overline{\rho}(D_j)) < 1$,
    daher $\rho^{\mathrm{sys}} = 1 - \prod_j(\ldots) > 0$.

    \textbf{(ii):} Sei $\overline{\rho}(D_k)$ wachsend. Dann f\"allt
    $(1 - w_k \overline{\rho}(D_k))$ und damit das gesamte Produkt. Also steigt
    $\rho^{\mathrm{sys}}$.

    \textbf{(iii):} $\rho^{\mathrm{sys}} < 1 \Leftrightarrow \prod_j(1-w_j\overline{\rho}(D_j)) > 0
    \Leftrightarrow \forall j\colon w_j\overline{\rho}(D_j) < 1
    \Leftrightarrow \overline{\rho}(D_j) < 1/w_j$. Da $w_j \leq 1$, ist
    $\overline{\rho}(D_j) < 1$ eine hinreichende Bedingung.

    \textbf{(iv):} Es gilt $1 - \prod_j(1-x_j) \geq x_k$ f\"ur alle $k$,
    wenn $x_j \in [0,1]$ (Bernoulli-Ungleichung in umgekehrter Form):
    $\prod_j(1-x_j) \leq (1-x_k)$, also $1 - \prod_j(1-x_j) \geq x_k$.
    Mit $x_j = w_j\overline{\rho}(D_j)$ folgt die Behauptung.
    $\qedhere$
\end{beweis}

\begin{korollar}[\"Ubereinstimmung mit dem Belastbarkeitsmodell]
\label{kor:uml-belastbarkeit}
    Das systemweite Risiko $\rho^{\mathrm{sys}}(\mathcal{D})$ ist konsistent mit
    dem Belastbarkeitsindex $B(t, u, r)$ aus Definition~\ref{def:belastbarkeit}:
    \[
        B(t, u, r) \;\leq\; 1 - \rho^{\mathrm{sys}}(\mathcal{D})
        \quad \text{f\"ur } t \geq t_0,
    \]
    wobei $t_0$ der Inbetriebnahmezeitpunkt ist. Das UML-Risikomodell liefert
    damit eine untere Schranke f\"ur den Belastbarkeitsindex ab Produktionsstart.
\end{korollar}

\begin{beweis}
    Nach Satz~\ref{satz:uml-konsistenz} gilt
    $B \leq 1 - \max_i \lambda_{e_i} \rho_{e_i}$. Nach Satz~\ref{satz:uml-aggregation}
    (iv) gilt $\rho^{\mathrm{sys}} \geq \max_j w_j \overline{\rho}(D_j) \geq
    \max_i \lambda_{e_i}\rho_{e_i} \cdot \min_j w_j$. F\"ur $\min_j w_j = 1/p$
    (Gleichgewichtung) folgt $\rho^{\mathrm{sys}} \geq \frac{1}{p}\max_i \lambda_{e_i}\rho_{e_i}$.
    Da $p$ endlich und $B \leq 1 - \max_i\lambda_{e_i}\rho_{e_i} \leq 1 - \rho^{\mathrm{sys}}$
    im dominanten Fall (ein Diagrammtyp tr\"agt das Maximum), folgt die Behauptung.
    $\qedhere$
\end{beweis}

% ---------------------------------------------------------------------------
\subsection{Relevanz der Diagrammtypen im Vergleich}
% ---------------------------------------------------------------------------

\begin{satz}[Diagrammtypspezifische Risikoschwergewichte]
\label{satz:diag-gewicht}
    Sei $\mathcal{S}$ ein Softwaresystem der Klasse~$k$ (gem\"a\ss{}
    Definition~\ref{def:klassen}). Die optimalen Diagrammgewichte $w_j^*$ f\"ur
    die Risikoaggregation erf\"ullen folgende Ungleichungskette:
    \begin{center}
    \renewcommand{\arraystretch}{1.2}
    \begin{tabular}{lcl}
        \toprule
        \textbf{Klasse} & \textbf{Dominanter Diagrammtyp} & \textbf{Begr\"undung} \\
        \midrule
        I   & Klassendiagramm & stabile Struktur dominiert \\
        II  & Komponentendiagramm & Deployment-Risiko dominiert \\
        III & Sequenzdiagramm & volatile Interaktionen dominieren \\
        IV  & alle gleichgewichtig & maximale Gesamtunsicherheit \\
        \bottomrule
    \end{tabular}
    \end{center}
\end{satz}

\begin{beweis}
    Die Begr\"undung folgt aus der Charakteristik der Vertragsklassen
    (Definition~\ref{def:klassen}) und den Risikodefinitionen der einzelnen
    Diagrammtypen:

    \textbf{Klasse~I} ($\kappa \leq 0.5$, $\nu \leq 0.5$): Geringe Komplexit\"at
    und geringe Volatili\-t\"at bedeuten, dass weder Interaktionen noch Deployment
    stark schwanken. Das dominante Risiko liegt in der statischen Struktur
    (Definition~\ref{def:strukturrisiko}), also im Klassendiagramm.

    \textbf{Klasse~II} ($\kappa > 0.5$, $\nu \leq 0.5$): Hohe technische Komplexit\"at
    bei stabilen Anforderungen f\"uhrt zu komplexen Deployment-Topologien
    (z.\,B.\ verteilte Systeme, Echtzeitsysteme). Das Komponentendiagramm mit
    \texttt{deployRisk} und Propagationsfaktoren dominiert.

    \textbf{Klasse~III} ($\kappa \leq 0.5$, $\nu > 0.5$): Volatile Anforderungen
    \"andern prim\"ar die Interaktionssequenzen zwischen Objekten. Nach
    Satz~\ref{satz:seq-superadditiv} akkumuliert sich das Risiko superadditiv in
    Sequenzdiagrammen; diese dominieren daher.

    \textbf{Klasse~IV} ($\kappa > 0.5$, $\nu > 0.5$): Maximale Unsicherheit in
    allen Dimensionen bedeutet, dass kein einzelner Diagrammtyp das Risiko allein
    repr\"asentieren kann. Gleichgewichtung $w_j = 1/p$ maximiert die Robustheit
    der Risikoabsch\"atzung gem\"a\ss{} Korollar~\ref{kor:uml-belastbarkeit}.
    $\qedhere$
\end{beweis}

% ===========================================================================
\section{Rechtskraft von Softwarevertr\"agen unter der Bedingung formaler\\
         Risikomodellierung}
\label{sec:rechtskraft}
% ===========================================================================

Die bisher entwickelten Ergebnisse dieser Arbeit --- das Belastbarkeitsmodell
$B(t,u,r)$, der Satz „Software hei\ss{}t Risiko" (Satz~\ref{satz:risiko}) und das
UML-Risikoprofil \textsf{RiskProfile} (Abschnitt~\ref{sec:uml}) --- f\"uhren zu einer
zentralen rechtlichen Konsequenz, die in diesem Kapitel formal ausgearbeitet wird:

\medskip
\begin{center}
\textit{Ein Softwarevertrag kann erst dann als rechtskr\"aftig gelten, wenn die\\
Risikomodellierung gr\"undlich und formal korrekt nachgewiesen wurde.}
\end{center}
\medskip

Diese These ist nicht nur eine praktische Empfehlung, sondern l\"asst sich aus den
formalen Ergebnissen dieser Arbeit logisch ableiten. Im Folgenden wird sie pr\"azisiert,
formal begr\"undet und in den rechtlichen Kontext eingebettet.

% ---------------------------------------------------------------------------
\subsection{Formale Vorbedingung der Vertragsg\"ultigkeit}
% ---------------------------------------------------------------------------

\begin{definition}[Risikonachweis]
\label{def:risikonachweis}
    Ein \emph{Risikonachweis} $\mathcal{N}$ f\"ur einen Softwarevertrag
    $\mathcal{V} = (\mathcal{S}, \Phi, t_0, T_{\max})$ ist ein Tupel:
    \[
        \mathcal{N} \;=\; \bigl(\Pi_R^*,\; \rho^{\mathrm{sys}*},\;
                               \mathcal{K}^*,\; \mathcal{P}^*\bigr),
    \]
    wobei
    \begin{itemize}
        \item $\Pi_R^*$: ein vollst\"andig annotiertes UML-Modell gem\"a\ss{}\
              Definition~\ref{def:riskprofile}, das alle OCL-Constraints $C_1$--$C_4$
              erf\"ullt,
        \item $\rho^{\mathrm{sys}*} \in (0,1)$: das berechnete systemweite Restrisiko
              gem\"a\ss{} Definition~\ref{def:systemrisiko}, formal hergeleitet aus
              $\Pi_R^*$,
        \item $\mathcal{K}^* \in \{\mathrm{I, II, III, IV}\}$: die zugewiesene
              Vertragsklasse gem\"a\ss{} Definition~\ref{def:klassen}, nachgewiesen
              durch Messung von $\kappa$ und $\nu$,
        \item $\mathcal{P}^*$: ein formaler Beweis oder eine maschinell gepr\"ufte
              Verifikation, dass $\Pi_R^*$ konsistent mit $\Phi$ und den Ergebnissen
              von Abschnitt~\ref{sec:uml} ist.
    \end{itemize}
\end{definition}

\begin{definition}[Rechtskr\"aftiger Softwarevertrag]
\label{def:rechtskraeftig}
    Ein Softwarevertrag $\mathcal{V}$ hei\ss{}t \emph{rechtskr\"aftig im Sinne dieser
    Arbeit}, wenn ein g\"ultiger Risikonachweis $\mathcal{N}$ im Sinne von
    Definition~\ref{def:risikonachweis} existiert und dieser Bestandteil des
    Vertragswerks ist:
    \[
        \text{rechtskr\"aftig}(\mathcal{V}) \;\Longleftrightarrow\;
        \exists\, \mathcal{N} \;\text{g\"ultig} \;\wedge\;
        \mathcal{N} \subseteq \mathcal{V}.
    \]
\end{definition}

% ---------------------------------------------------------------------------
\subsection{Notwendigkeit des Riskonachweises: Formale Begr\"undung}
% ---------------------------------------------------------------------------

\begin{satz}[Notwendigkeit des Risikona\-chweis\-es f\"ur Vertragsg\"ultigkeit]
\label{satz:notwendigkeit}
    Sei $\mathcal{V}$ ein Softwarevertrag ohne g\"ultigen Risikonachweis $\mathcal{N}$.
    Dann existiert stets eine Situation, in der eine Vertragspartei eine
    Verpflichtung eingeht, deren Erf\"ullbarkeit strukturell ungewiss ist ---
    was die informierte Einwilligung (\emph{informed consent}) beider Parteien
    unm\"oglich macht.
\end{satz}

\begin{beweis}
    Der Beweis verl\"auft in drei Schritten.

    \textbf{Schritt~1: Strukturelles Restrisiko ohne Nachweis.}
    Nach Satz~\ref{satz:risiko} gilt f\"ur jedes nicht-triviale Softwaresystem
    $\mathcal{S}$:
    \[
        P\bigl(\exists\, z \in \mathcal{Z}(\mathcal{S}) \setminus \mathcal{T}
        \colon \mathcal{S}(z,E) \notin \Phi\bigr) \;>\; 0.
    \]
    Ohne Risikonachweis $\mathcal{N}$ ist dieses Risiko den Vertragsparteien weder
    quantifiziert noch kommuniziert. Es existiert also ein positives, den Parteien
    unbekanntes Restrisiko.

    \textbf{Schritt~2: Fehlende informierte Einwilligung.}
    Ein Vertrag im Sinne des BGB (\S~145~ff.) setzt die Willenserkl\"arung beider
    Parteien voraus. Eine Willenserkl\"arung ist nur dann inhaltlich bestimmt, wenn
    der Gegenstand der Verpflichtung hinreichend klar ist \cite{kilian2017}. Da ohne
    $\mathcal{N}$ das systemweite Risiko $\rho^{\mathrm{sys}}$ unbekannt ist, kann
    der Auftraggeber nicht beurteilen, ob er die mit dem Softwaresystem verbundenen
    Risiken akzeptiert --- seine Willenserkl\"arung ist informationsdefizit\"ar.

    \textbf{Schritt~3: Logische \"Aquivalenz.}
    Sei $I(P)$ die Aussage „Partei $P$ hat informierte Kenntnis der
    Vertragsverpflichtung". Ohne $\mathcal{N}$ gilt:
    $\neg I(\text{Auftraggeber})$, denn $\rho^{\mathrm{sys}}$ ist unbekannt.
    Da ein g\"ultiger Vertrag $I(\text{Auftraggeber}) \wedge
    I(\text{Auftragnehmer})$ voraussetzt \cite{kilian2017}, folgt:
    $\neg \mathcal{N} \Rightarrow \neg \text{rechtskr\"aftig}(\mathcal{V})$.
    Kontraposition: $\text{rechtskr\"aftig}(\mathcal{V}) \Rightarrow
    \exists\, \mathcal{N}$. $\qedhere$
\end{beweis}

\begin{korollar}[Risikonachweis als Vertragsbedingung]
\label{kor:bedingung}
    Der Risikonachweis $\mathcal{N}$ ist eine \emph{notwendige Bedingung} f\"ur die
    Rechtskraft eines Softwarevertrags. Kein Softwarevertrag, der das strukturelle
    Restrisiko des zu entwickelnden Systems verschweigt oder nicht formal quantifiziert,
    kann als vollst\"andig informierter und damit rechtskr\"aftiger Vertrag gelten.
\end{korollar}

% ---------------------------------------------------------------------------
\subsection{Hinreichende Bedingungen: Vollst\"andigkeit des Risikona\-chweis\-es}
% ---------------------------------------------------------------------------

\begin{definition}[Vollst\"andigkeit des Risikona\-chweis\-es]
\label{def:vollstaendigkeit-rn}
    Ein Risikonachweis $\mathcal{N}$ hei\ss{}t \emph{vollst\"andig}, wenn er
    folgende vier Bedingungen erf\"ullt:
    \begin{enumerate}[label=(\roman*)]
        \item \textbf{Modellvollst\"andigkeit}: Das annotierte UML-Modell $\Pi_R^*$
              umfasst mindestens alle Diagrammtypen, die nach Satz~\ref{satz:diag-gewicht}
              f\"ur die Vertragsklasse $\mathcal{K}^*$ dominant sind.
        \item \textbf{Quantitative Nachvollziehbarkeit}: Alle Tagged Values
              $(\rho_e, u_e, \lambda_e, \kappa_e)$ in $\Pi_R^*$ sind entweder durch
              Messung (statische Analyse, Coverage-Metriken) oder durch begr\"undete
              Sch\"atzung mit angegebenem Konfidenzintervall belegt.
        \item \textbf{Aggregationskonsistenz}: Das berechnete systemweite Restrisiko
              $\rho^{\mathrm{sys}*}$ folgt nachvollziehbar aus den Elementrisiken
              gem\"a\ss{} der Formel in Definition~\ref{def:systemrisiko}.
        \item \textbf{Risikoklauseln}: Der Vertrag $\mathcal{V}$ enth\"alt explizite
              Risikoklauseln $\mathcal{R}_C$ gem\"a\ss{} Definition~\ref{def:risikonachweis},
              die $\rho^{\mathrm{sys}*}$ benennen und Haftung sowie Mitigationsma\ss{}nahmen
              regeln.
    \end{enumerate}
\end{definition}

\begin{satz}[Hinreichende Bedingung f\"ur Rechtskraft]
\label{satz:hinreichend}
    Sei $\mathcal{N}$ ein vollst\"andiger Risikonachweis im Sinne von
    Definition~\ref{def:vollstaendigkeit-rn}. Dann sind beide Vertragsparteien
    in der Lage, eine informierte Willenserkl\"arung abzugeben, und der Vertrag
    $\mathcal{V}$ erf\"ullt die formalen Voraussetzungen der Rechtskraft.
\end{satz}

\begin{beweis}
    Aus Bedingung~(i) folgt, dass alle risikodominanten Systemaspekte modelliert
    sind. Aus Bedingung~(ii) folgt, dass die Risikoma\ss{}zahlen nicht beliebig
    angesetzt, sondern empirisch begr\"undet sind --- der Auftraggeber kann sie
    nachpr\"ufen oder durch eigene Sachverst\"andige pr\"ufen lassen. Aus
    Bedingung~(iii) folgt, dass das Gesamtrisiko $\rho^{\mathrm{sys}*}$ formal
    aus den Elementen ableitbar ist und keine willk\"urliche Gesamtangabe darstellt.
    Aus Bedingung~(iv) folgt, dass das Restrisiko nicht nur berechnet, sondern
    auch vertraglich adressiert ist.

    Damit gilt: Beide Parteien kennen Art und Gr\"o\ss{}e des Risikos, seine
    Verteilung auf Systemkomponenten und die vereinbarte Reaktion im Schadensfall.
    Die Voraussetzung der informierten Willenserkl\"arung aus Satz~\ref{satz:notwendigkeit}
    ist erf\"ullt. Die formalen Anforderungen an einen bestimmten und einklagbaren
    Vertragsinhalt nach \S\S~145,\,241~BGB sind gegeben. $\qedhere$
\end{beweis}

% ---------------------------------------------------------------------------
\subsection{Stufenmodell der Vertragsg\"ultigkeit}
% ---------------------------------------------------------------------------

Aus den vorstehenden Definitionen und S\"atzen ergibt sich ein dreistufiges
Modell der Vertragsg\"ultigkeit, das den Zusammenhang zwischen
Risikomodellierungstiefe und Rechtskraft strukturiert:

\begin{definition}[Stufenmodell der Softwarevertrags-Rechtskraft]
\label{def:stufenmodell}
    Sei $\mathcal{V}$ ein Softwarevertrag. Wir definieren drei G\"ultigkeitsstufen:
    \begin{center}
    \renewcommand{\arraystretch}{1.3}
    \begin{tabular}{clp{7.2cm}}
        \toprule
        \textbf{Stufe} & \textbf{Bezeichnung} & \textbf{Bedingung} \\
        \midrule
        0 & Ung\"ultig
          & Kein Risikonachweis; strukturelles Restrisiko unbekannt.
            Nach Satz~\ref{satz:notwendigkeit} keine informierte Einwilligung
            m\"oglich. \\[4pt]
        1 & Bedingt g\"ultig
          & Risikonachweis $\mathcal{N}$ vorhanden, aber unvollst\"andig
            (nicht alle Bedingungen von Definition~\ref{def:vollstaendigkeit-rn}
            erf\"ullt). Vertrag ist anfechtbar, da Informationsdefizite bestehen.\\[4pt]
        2 & Vollst\"andig rechtskr\"aftig
          & Vollst\"andiger Risikonachweis $\mathcal{N}$ gem\"a\ss{}
            Definition~\ref{def:vollstaendigkeit-rn}. Nach Satz~\ref{satz:hinreichend}
            sind alle Voraussetzungen informierter Willenserkl\"arung erf\"ullt. \\
        \bottomrule
    \end{tabular}
    \end{center}
\end{definition}

\begin{lemma}[Monotonie der G\"ultigkeitsstufen]
\label{lem:stufen-monoton}
    Sei $\mathcal{N}_1$ ein Risikonachweis, der $k$ der vier Bedingungen aus
    Definition~\ref{def:vollstaendigkeit-rn} erf\"ullt, und $\mathcal{N}_2$ ein
    Risikonachweis, der $k' > k$ Bedingungen erf\"ullt. Dann gilt:
    \[
        \mathrm{G\"ultigkeitsstufe}(\mathcal{N}_2) \;\geq\;
        \mathrm{G\"ultigkeitsstufe}(\mathcal{N}_1).
    \]
    Mehr erf\"ullte Nachweisbedingungen erh\"ohen die Vertragsg\"ultigkeit
    monoton.
\end{lemma}

\begin{beweis}
    Da jede zus\"atzlich erf\"ullte Bedingung ein bestehendes Informationsdefizit
    schlie\ss{}t (Modellvollst\"andigkeit, Quantifizierbarkeit, Aggregationskonsistenz
    oder vertragliche Verankerung), n\"ahert sich der Risikonachweis monoton der
    vollst\"andigen informierten Einwilligung. Die Stufenzuweisung in
    Definition~\ref{def:stufenmodell} ist eine Treppenfunktion in $k$, die bei
    $k=4$ Stufe~2 erreicht. Da $k' > k$, gilt $\mathrm{Stufe}(k') \geq \mathrm{Stufe}(k)$.
    $\qedhere$
\end{beweis}

% ---------------------------------------------------------------------------
\subsection{Konsequenzen f\"ur die Vertragsgestaltung in der Praxis}
% ---------------------------------------------------------------------------

\begin{bemerkung}[Pr\"aufpflicht und Due Diligence]
    Aus Korollar~\ref{kor:bedingung} und Satz~\ref{satz:hinreichend} ergibt sich f\"ur
    die Vertragspraxis eine klare \emph{Due-Diligence-Pflicht}: Auftraggeber und
    Auftragnehmer sind verpflichtet, vor Vertragsschluss gemeinsam einen vollst\"andigen
    Risikonachweis $\mathcal{N}$ zu erstellen. Diese Pflicht ist nicht delegierbar ---
    sie erfordert, dass beide Parteien die in $\Pi_R^*$ enthaltenen Risikoma\ss{}zahlen
    verstehen und akzeptieren. In der Praxis empfiehlt sich die Einbindung eines
    unabh\"angigen Sachverst\"andigen, der den Risikonachweis pr\"uft und best\"atigt.
\end{bemerkung}

\begin{bemerkung}[Vertragsklasse und Nachweistiefe]
    Die erforderliche Nachweistiefe ist abh\"angig von der Vertragsklasse
    $\mathcal{K}^*$. F\"ur Klasse~I (geringe Komplexit\"at, stabile Anforderungen)
    kann ein vereinfachter Nachweis gen\"ugen; f\"ur Klasse~IV (hohe Komplexit\"at,
    volatile Anforderungen) ist nach Satz~\ref{satz:diag-gewicht} eine vollst\"andige
    Annotation aller f\"unf Diagrammtypen unabdingbar. Dies korrespondiert mit
    Satz~\ref{satz:solo}: Gerade bei Klasse~IV ben\"otigt der Risikonachweis selbst
    ein interdisziplin\"ares Team.
\end{bemerkung}

\begin{hypothese}[Geringere Streitigkeitsrate bei vollst\"andigem Nachweis]
\label{hyp:streit}
    Softwarevertr\"age, die einen vollst\"andigen Risikonachweis $\mathcal{N}$
    gem\"a\ss{} Definition~\ref{def:vollstaendigkeit-rn} enthalten, weisen eine
    statistisch signifikant geringere Rate an Vertragsstreitigkeiten auf als Vertr\"age
    ohne Risikonachweis, da Erwartungen und Risiken ex ante klar kommuniziert sind.
\end{hypothese}

\begin{bemerkung}[Verhältnis zu bestehenden Normen]
    Das hier entwickelte Stufenmodell ist kompatibel mit und baut auf bestehenden
    Normen auf. ISO/IEC~25010 definiert Qualit\"atsmerkmale, die in die Tagged Values
    des \textsf{RiskProfile} einflie\ss{}en. ISO/IEC~12207 beschreibt den
    Softwarelebenszyklusprozess, in den der Risikonachweis als eigene Phase
    „Risikomodellierung und -nachweis" integriert werden kann --- idealerweise zwischen
    Anforderungsanalyse und Vertragsschluss. IEC~61508 fordert f\"ur sicherheitskritische
    Systeme bereits eine systematische Risikoanalyse; das vorliegende Stufenmodell
    verallgemeinert diesen Gedanken auf \emph{alle} Softwarevertr\"age, unabh\"angig
    von der Dom\"ane.
\end{bemerkung}

% ===========================================================================
\section{Zusammenfassung}
% ===========================================================================

Diese Arbeit hat ein umfassendes formales Rahmenwerk zur Belastbarkeit von
Softwarevertr\"agen, zur quantitativen Risikomodellierung in UML und zur
rechtlichen G\"ultigkeit von Softwarevertr\"agen entwickelt. Die zentralen
Ergebnisse lassen sich wie folgt zusammenfassen:

\begin{itemize}
    \item Softwaresysteme k\"onnen aufgrund des Halteproblems und der unendlichen
          Zustandsmenge nicht vollst\"andig formal verifiziert werden
          (Satz~\ref{satz:unvollstaendigkeit}, Satz~\ref{satz:tests}).

    \item Der Belastbarkeitsindex $B(t,u,r)$ f\"allt monoton mit Zeit, Unsicherheit
          und Benutzerzahl (Satz~\ref{satz:verfall}). Kein Softwarevertrag bleibt
          dauerhaft belastbar (Korollar~\ref{kor:langfristig}).

    \item Das Restrisiko nach Testabschluss ist strukturell positiv und nicht
          eliminierbar (Korollar~\ref{kor:restrisiko}).

    \item \textbf{Software hei\ss{}t Risiko}: Absolute Produktionssicherheit ist
          nicht garantierbar (Satz~\ref{satz:risiko}). Dieses Risiko ist normal
          und muss von allen Beteiligten bewusst eingeplant werden.

    \item Vier Vertragsklassen (I--IV) erm\"oglichen eine systematische
          Risikokategorisierung nach technischer Komplexit\"at und
          Anforderungsvolatili\-t\"at (Definition~\ref{def:klassen},
          Satz~\ref{satz:diag-gewicht}).

    \item F\"ur Klassen~III und~IV ist ein interdisziplin\"ares Team zwingend
          erforderlich (Satz~\ref{satz:solo}).

    \item Das UML-Risikoprofil \textsf{RiskProfile} erm\"oglicht die quantitative
          Annotation von Risikoma\ss{}zahlen direkt in UML-Struktur- und
          Verhaltensdiagrammen (Abschnitt~\ref{sec:uml}). Die Risikoannotation ist
          formal konsistent mit dem Belastbarkeitsmodell $B(t,u,r)$ und erlaubt
          eine diagrammtypspezifische Risikoaggregation
          (S\"atze~\ref{satz:uml-monoton}--\ref{satz:uml-aggregation}).

    \item \textbf{Rechtskraft setzt Risikonachweis voraus}: Ein Softwarevertrag
          kann erst dann als rechtskr\"aftig gelten, wenn ein vollst\"andiger
          Risikonachweis $\mathcal{N}$ gem\"a\ss{} Definition~\ref{def:vollstaendigkeit-rn}
          vorliegt und Bestandteil des Vertragswerks ist
          (Satz~\ref{satz:notwendigkeit}, Korollar~\ref{kor:bedingung},
          Satz~\ref{satz:hinreichend}). Das dreistufige G\"ultigkeitsmodell
          (Definition~\ref{def:stufenmodell}) ordnet Vertr\"age nach ihrer
          Risikomodellierungstiefe.
\end{itemize}

% ===========================================================================
\section{Ausblick}
% ===========================================================================

Die vorliegende Arbeit er\"offnet eine Vielzahl von Forschungsrichtungen. Der Ausblick
gliedert sich in drei Teile: Zun\"achst werden Perspektiven aus dem Kapitel zur
Rechtskraft von Softwarevertr\"agen (Abschnitt~\ref{sec:rechtskraft}) entwickelt,
anschlie\ss{}end folgen Perspektiven aus dem UML-Risikokapitel
(Abschnitt~\ref{sec:uml}), und abschlie\ss{}end die \"ubergreifenden
Forschungsrichtungen zum Belastbarkeitsmodell und seinem weiteren wissenschaftlichen
Umfeld.

% ---------------------------------------------------------------------------
\subsection{Empirische Validierung des Stufenmodells der Vertragsg\"ultigkeit}
% ---------------------------------------------------------------------------

Das in Definition~\ref{def:stufenmodell} entwickelte dreistufige G\"ultigkeitsmodell
ist ein formales Konstrukt, das bisher nicht empirisch validiert wurde.
Hypothese~\ref{hyp:streit} besagt, dass Vertr\"age mit vollst\"andigem Risikonachweis
eine geringere Streitigkeitsrate aufweisen. Eine empirische Untersuchung k\"onnte
diese Hypothese anhand von Urteilsdatenbanken (z.\,B.\ juris, beck-online,
ECLI-Datenbanken) testen: Wie h\"aufig wurde in Softwarevertragsstreitigkeiten das
Fehlen einer dokumentierten Risikoabsch\"atzung als Ursache identifiziert? Welche
Schadensh\"ohen sind typisch f\"ur Vertr\"age der Klassen~I--IV? Diese Daten k\"onnten
die Gewichte des Stufenmodells kalibrieren und seine rechtspraktische Relevanz belegen.

% ---------------------------------------------------------------------------
\subsection{Normative Verankerung des Risikona\-chweis\-es im Vertragsrecht}
% ---------------------------------------------------------------------------

Satz~\ref{satz:notwendigkeit} und Korollar~\ref{kor:bedingung} begr\"unden den
Risikonachweis als notwendige Bedingung f\"ur informierte Willenserkl\"arungen.
Diese Erkenntnis ist bisher nicht in positivem Recht kodifiziert. Zuk\"unftige
Arbeiten k\"onnten konkrete Gesetzgebungsvorschl\"age entwickeln:
\begin{itemize}
    \item Eine \textbf{Erg\"anzung von \S~631 BGB} (Werkvertragsrecht) um eine
          explizite Risikooffenbarungspflicht f\"ur Softwareentwicklungsvertr\"age
          ab einer definierten Komplexit\"atsschwelle $\kappa > 0.5$ (Klassen~II
          und~IV).
    \item Eine \textbf{EU-Richtlinie zur Softwarevertragstransparenz}, die ---
          analog zur Verbraucherrechterichtlinie 2011/83/EU --- vorschreibt,
          dass dem Auftraggeber vor Vertragsschluss ein verst\"andlicher
          Risikonachweis \"ubergeben werden muss.
    \item Die \textbf{Integration in den EU-AI-Act} \cite{euai2024}: F\"ur
          KI-basierte Hochrisikosysteme (Anhang~III AI~Act) k\"onnte das
          vollst\"andig annotierte \textsf{RiskProfile} als Teil der technischen
          Dokumentation nach Anhang~IV gefordert werden.
\end{itemize}

% ---------------------------------------------------------------------------
\subsection{Schiedsgerichtliche Nutzung des Risikona\-chweis\-es}
% ---------------------------------------------------------------------------

Im Streitfall bietet der formale Risikonachweis $\mathcal{N}$ eine sachliche
Grundlage f\"ur Schiedsverfahren: Da $\rho^{\mathrm{sys}*}$ und die diagrammbasierten
Elementrisiken $\omega(e_i)$ numerisch vorliegen, kann ein technischer Sachverst\"andiger
pr\"ufen, ob ein aufgetretener Fehler im vertraglich dokumentierten Risikobereich lag
oder einen unvorhergesehenen Ausnahmefall darstellt. Dies w\"urde eine objektivere
Haftungsabgrenzung erm\"oglichen: Liegt der Fehler in einem Element mit
$\lambda_e \cdot \rho_e > \rho^{\mathrm{sys}*}$, war er im Risikonachweis
antizipiert; liegt er au\ss{}erhalb, ist eine Nachtragsregelung sachgerecht.
Zuk\"unftige Arbeiten k\"onnten Musterschiedsklauseln entwickeln, die den
Risikonachweis als Beweismittel in Softwarevertragsstreitigkeiten etablieren.

% ---------------------------------------------------------------------------
\subsection{Werkzeugunterst\"utzung f\"ur das \textsf{RiskProfile}}
% ---------------------------------------------------------------------------

Das in Definition~\ref{def:riskprofile} definierte UML-Profil \textsf{RiskProfile}
ist derzeit ein rein formales Konstrukt. Ein wesentlicher n\"achster Schritt ist seine
\emph{Implementierung als maschinenlesbare Profilerdatei} im XMI-Format (XML Metadata
Interchange), dem OMG-Standard f\"ur den Austausch von UML-Modellen \cite{omg2017}.
Konkret sind folgende Werkzeugintegrationen denkbar:

\begin{itemize}
    \item \textbf{Eclipse Papyrus}: Das quelloffene UML-Modellierungswerkzeug
          unterst\"utzt Profilerdateien nativ. Eine \textsf{RiskProfile}-Erweiterung
          k\"onnte als Eclipse-Plugin verteilt werden und die Tagged-Value-Eingabe
          f\"ur alle sechs Stereotypen aus Definition~\ref{def:riskprofile} mit
          Eingabevalidierung gegen\"uber den OCL-Constraints $C_1$--$C_4$ erm\"oglichen.
    \item \textbf{Enterprise Architect (Sparx Systems)}: Weit verbreitetes kommerzielles
          Werkzeug mit MDG-Technologie (Model Driven Generation), \"uber die
          \textsf{RiskProfile} als angepasstes Technologiepaket eingebettet werden kann.
    \item \textbf{PlantUML-Erweiterung}: F\"ur die leichtgewichtige,
          textbasierte UML-Erstellung k\"onnte eine DSL-Erweiterung definiert werden,
          die Risikoannotationen in der Syntax \texttt{@risk(0.82, 0.9)} direkt in
          Diagrammquelltext erm\"oglicht und automatisch in Risikosummaries \"ubersetzt.
\end{itemize}

Eine solche Implementierung w\"urde es Projektteams erstmals erlauben, die in
Satz~\ref{satz:uml-aggregation} definierte systemweite Risikoaggregation
$\rho^{\mathrm{sys}}(\mathcal{D})$ automatisiert und diagramm\"ubergreifend zu berechnen.

% ---------------------------------------------------------------------------
\subsection{Automatische Risikoberechnung aus Quellcode und Modell}
% ---------------------------------------------------------------------------

Die Risikoma\ss{}zahlen $\omega(e) = (\rho_e, u_e, \lambda_e, \kappa_e)$ aus
Definition~\ref{def:risikozahl} werden in der vorliegenden Arbeit manuell annotiert.
Eine wichtige Forschungsrichtung ist die \emph{automatische Extraktion} dieser Werte
aus bestehenden Artefakten:

\begin{itemize}
    \item \textbf{Kopplungskomplexit\"at $\kappa_e$}: Fan-in- und Fan-out-Metriken
          lassen sich aus Java-Bytecode, Python-AST-Analysen oder C++-Include-Graphen
          automatisch berechnen. Werkzeuge wie SonarQube oder Understand liefern
          bereits entsprechende Metriken, die in \texttt{coupling}-Tagged-Values
          \"ubertragen werden k\"onnen.
    \item \textbf{Zustandsraumsch\"atzung f\"ur $\rho^{\mathrm{state}}$}: Symbolische
          Ausf\"uhrung (z.\,B.\ mit KLEE oder Java PathFinder) kann den erreichbaren
          Zustandsraum $\mathcal{Z}(\mathcal{S})$ approximieren und damit den
          Testabdeckungsgrad aus Definition~\ref{def:zustandsrisiko} quantifizieren.
    \item \textbf{Zyklomatische Komplexit\"at f\"ur $\rho^{\mathrm{ctrl}}$}: McCabes
          Ma\ss{} $M(G_A)$ aus Definition~\ref{def:kfrisiko} ist durch statische
          Codeanalyse direkt berechenbar und kann automatisch in die
          \texttt{residualRisk}-Tagged-Values von Aktivit\"atsdiagrammelementen
          \"ubertragen werden.
\end{itemize}

Ein vollst\"andiger \emph{Risk-Annotation-Pipeline}-Ansatz --- von Quellcode und
vorhandenen UML-Modellen bis hin zu einem vollst\"andig annotierten
\textsf{RiskProfile}-Modell --- w\"are ein bedeutender Beitrag zur
modellgetriebenen Qualit\"atssicherung.

% ---------------------------------------------------------------------------
\subsection{Validierung des Risikoaggregationstheorems an realen Projekten}
% ---------------------------------------------------------------------------

Satz~\ref{satz:uml-aggregation} (Risikoaggregationstheorem) und
Korollar~\ref{kor:uml-belastbarkeit} (\"Ubereinstimmung mit dem Belastbarkeitsmodell)
stellen formale Resultate dar, die bisher nicht empirisch validiert wurden. Eine
wesentliche Forschungsaufgabe besteht darin, anhand realer Projektdaten zu pr\"ufen,
ob $\rho^{\mathrm{sys}}(\mathcal{D})$ tats\"achlich eine g\"ultige untere Schranke
f\"ur $1 - B(t, u, r)$ liefert. Hierf\"ur bieten sich Open-Source-Projekte mit
langer Versionierungsgeschichte an (z.\,B.\ Linux-Kernel, Apache Commons, Eclipse
IDE), bei denen sowohl UML-Modelle als auch Bug-Datenbanken und Ausfallhistorien
vorliegen. Ein Vergleich zwischen den vorhergesagten Risikowerten aus
$\rho^{\mathrm{sys}}$ und den tats\"achlich beobachteten Fehlerh\"aufigkeiten w\"urde
die Vorhersagequalit\"at des Modells messen und ggf.\ eine Rekalibrierung der
Diagrammgewichte $w_j$ aus Satz~\ref{satz:diag-gewicht} motivieren.

% ---------------------------------------------------------------------------
\subsection{Erweiterung des \textsf{RiskProfile} um zeitliche Dynamik}
% ---------------------------------------------------------------------------

Das aktuelle \textsf{RiskProfile} annotiert statische Risikoma\ss{}zahlen, die zum
Modellierungszeitpunkt gelten. Da der Belastbarkeitsindex $B(t, u, r)$ aber ein
zeitdynamisches Modell ist (Satz~\ref{satz:verfall}), liegt es nahe, auch die
UML-Risikoannotationen mit Zeitparametern zu versehen. Ein erweitertes Tagged Value
\texttt{riskDecayRate: Real} k\"onnte f\"ur jeden Stereotyp die Rate modellieren,
mit der das Restrisiko des Elements mit der Zeit w\"achst:
\[
    \rho_e(t) \;=\; 1 - e^{-\alpha_e \cdot t}, \quad \alpha_e > 0.
\]
Damit w\"are es m\"oglich, zeitliche Risikoentwicklungen direkt im UML-Modell
sichtbar zu machen --- etwa durch eine farbkodierte Animation in einem
UML-Modellierungswerkzeug, die zeigt, wie einzelne Klassen, Komponenten und
Nachrichten mit wachsender Systemlaufzeit riskanter werden.

% ---------------------------------------------------------------------------
\subsection{Integration des \textsf{RiskProfile} in agile Entwicklungsprozesse}
% ---------------------------------------------------------------------------

Agile Methoden wie Scrum oder Kanban planen in kurzen Iterationen und reagieren
flexibel auf ver\"anderte Anforderungen --- genau jene Situation, die Klasse~III
und~IV in Definition~\ref{def:klassen} beschreibt. Das \textsf{RiskProfile} bietet
hier einen formalen Ankerpunkt: Die Risikoma\ss{}zahlen k\"onnen als
\emph{Definition of Ready} und \emph{Definition of Done} in User Stories integriert
werden. Konkret k\"onnte ein Sprint-Abschluss verlangen, dass kein Element mit
$\lambda_e \cdot \rho_e > 0.5$ ohne Risikobehandlungsma\ss{}nahme in den
Produktionsbranch \"ubernommen wird. Dieses Kriterium ist direkt aus
Satz~\ref{satz:uml-konsistenz} ableitbar und w\"urde die formalen Ergebnisse dieser
Arbeit in die t\"agliche Entwicklungspraxis \"ubertragen.

% ---------------------------------------------------------------------------
\subsection{\textsf{RiskProfile} und modellbasierte Sicherheitsnachweise}
% ---------------------------------------------------------------------------

In sicherheitskritischen Dom\"anen (Automotive nach ISO~26262, Avionik nach
DO-178C, Medizintechnik nach IEC~62304) ist der Nachweis der funktionalen Sicherheit
Pflicht. Diese Normen fordern u.\,a.\ eine systematische Risikoanalyse, die bislang
in getrennten Artefakten (FMEA, FTA, HAZOP) gef\"uhrt wird. Das \textsf{RiskProfile}
k\"onnte hier als \emph{Br\"ucke zwischen UML-Systemmodell und Sicherheitsnachweis}
dienen: Indem die Risikoma\ss{}zahlen $\omega(e)$ direkt im Systemmodell annotiert
sind, entf\"allt die manuelle \"Ubertragung zwischen Modell und Sicherheitsartefakt.
Insbesondere k\"onnte das Propagationsmodell aus Satz~\ref{satz:komp-propagation}
als formale Grundlage f\"ur eine modellbasierte FMEA-Automatisierung dienen: Aus
dem Komponentenrisikograph $\mathcal{G} = (V, E_K)$ lassen sich Fehlerbaumstrukturen
algorithmisch ableiten, was den Nachweisaufwand erheblich reduzieren w\"urde.

% ---------------------------------------------------------------------------
\subsection{Diagrammtypgewichte und maschinelles Lernen}
% ---------------------------------------------------------------------------

Satz~\ref{satz:diag-gewicht} legt eine qualitative Begr\"undung f\"ur die optimalen
Diagrammgewichte $w_j^*$ je Vertragsklasse vor. Eine datengetriebene Alternative
besteht darin, diese Gewichte durch maschinelles Lernen aus historischen Projektdaten
zu sch\"atzen. Gegeben eine Datenbank von abgeschlossenen Projekten mit bekannter
Vertragsklasse, gemessenen Fehlerh\"aufigkeiten pro Diagrammtyp und beobachteten
Belastbarkeitsverl\"aufen $B(t)$, k\"onnten Regressionsmodelle oder neuronale Netze
die Parameter $w_j$ so kalibrieren, dass die Vorhersage von $\rho^{\mathrm{sys}}$
den beobachteten Fehlerquoten am besten entspricht. Ans\"atze wie
\emph{Survival Analysis} mit zeit\"abhangigen Kovariaten \cite{cox1972} bieten
hierf\"ur einen nat\"urlichen statistischen Rahmen.

% ---------------------------------------------------------------------------
\subsection{Quantitative Kalibrierung des Belastbarkeitsmodells}
% ---------------------------------------------------------------------------

Das in Definition~\ref{def:belastbarkeit} eingef\"uhrte Modell $B(t, u, r) = e^{-\delta(u,r)\,t}$
enth\"alt die Parameter $\delta_0, \alpha, \beta$, die in dieser Arbeit als qualitativ
gegeben behandelt wurden. Eine naheliegende Forschungsaufgabe ist die \emph{empirische
Kalibrierung} dieser Parameter anhand realer Projektdaten. Hierf\"ur k\"onnen Datenbanken
wie der ISBSG-Datensatz \cite{isbsg2020}, die Chaos-Reports der Standish-Group \cite{standish2020}
oder interne Projektarchive von Softwareunternehmen verwendet werden. Regressionsanalysen
und maschinelles Lernen k\"onnen dabei helfen, dom\"anenspezifische Parameterwerte zu
sch\"atzen. Die in Abschnitt~\ref{sec:uml} hergeleiteten Risikoma\ss{}zahlen
$\rho^{\mathrm{struct}}$, $\rho^{\mathrm{seq}}$ und $\rho^{\mathrm{prop}}$ bieten
dabei zus\"atzliche erkl\"arende Variablen, die die Kalibrierungsqualit\"at gegen\"uber
rein zeitbasierten Modellen verbessern k\"onnten.

% ---------------------------------------------------------------------------
\subsection{Stochastische Erweiterung}
% ---------------------------------------------------------------------------

Das aktuelle Modell ist deterministisch. Eine wichtige Erweiterung besteht darin,
$B(t, u, r)$ als stochastischen Prozess zu formulieren. Denkbar ist ein
\emph{geometrischer Brownscher Prozess}:
\[
    dB = -\delta(u,r)\, B\, dt + \sigma\, B\, dW_t,
\]
wobei $W_t$ ein Wiener-Prozess und $\sigma > 0$ die Volatilit\"at des Projekts ist.
Dieses Modell w\"urde es erm\"oglichen, Konfidenzintervalle f\"ur $B(t)$ zu berechnen
und das Risiko eines pl\"otzlichen Vertragsversagens zu quantifizieren. Die
stochastischen Risikoma\ss{}zahlen aus dem \textsf{RiskProfile} ---
insbesondere \texttt{failProb} und \texttt{propagationFactor} --- k\"onnten dabei als
diagrammbasierte Sch\"atzer f\"ur $\sigma$ dienen und die stochastische Modellierung
empirisch verankern.

% ---------------------------------------------------------------------------
\subsection{Spieltheoretische Modellierung}
% ---------------------------------------------------------------------------

Softwarevertr\"age involvieren mindestens zwei strategisch handelnde Akteure: Auftraggeber
und Auftragnehmer. Beide haben asymmetrische Information: Der Auftragnehmer kennt die
technischen Risiken besser; der Auftraggeber kennt seine gesch\"aftlichen Anforderungen
besser. Diese Situation l\"asst sich als \emph{Bayesianisches Spiel} modellieren
\cite{harsanyi1967}, in dem Gleichgewichte der Informationsweitergabe (\emph{cheap talk},
\emph{signaling}) untersucht werden. Wann ist es f\"ur den Auftragnehmer rational, Risiken
offenzulegen? Das \textsf{RiskProfile} k\"onnte hier als \emph{Signaling-Instrument}
dienen: Ein Auftragnehmer, der sein UML-Modell vollst\"andig mit \textsf{RiskProfile}
annotiert und dem Auftraggeber zur Verf\"ugung stellt, sendet ein glaubw\"urdiges Signal
\"uber seine Risikobereitschaft zur Transparenz --- was nach Mechanism-Design-Argumenten
\cite{harsanyi1967} vertrauensbildend und vertragsst\"arkend wirkt.

% ---------------------------------------------------------------------------
\subsection{Formale Methoden und Vertragsprüfung}
% ---------------------------------------------------------------------------

Die Ergebnisse dieser Arbeit legen nahe, dass formale Methoden --- wie Model Checking
\cite{clarke1999}, abstrakte Interpretation \cite{cousot1977} oder Hoare-Logik
\cite{hoare1969} --- systematischer in den Vertragslebenszyklus eingebettet werden
sollten. Konkret k\"onnten formale Spezifikationssprachen wie TLA+ \cite{lamport2002},
Alloy oder Event-B dazu verwendet werden, kritische Teile der Vertragsspezifikation
$\Phi$ maschinenverifikabel zu machen. Eine besonders vielversprechende Verkn\"upfung
besteht zwischen dem \textsf{RiskProfile} und Alloy: Die OCL-Constraints
$C_1$--$C_4$ aus Definition~\ref{def:riskprofile} k\"onnten direkt in Alloy-Pr\"adikate
\"ubersetzt werden, sodass die strukturelle Konsistenz des annotierten Modells automatisch
gepr\"uft werden kann. Dies w\"urde zwar nicht das Restrisiko aus Korollar~\ref{kor:restrisiko}
eliminieren, aber f\"ur die spezifizierten Teile eine nachweislich h\"ohere Belastbarkeit
erreichen.

% ---------------------------------------------------------------------------
\subsection{Rechtsinformatik: Softwarevertr\"age und Produkthaftung}
% ---------------------------------------------------------------------------

Die Frage der Haftung f\"ur Softwarefehler im Produktionsbetrieb ist in der Rechtspraxis
nach wie vor ungekl\"art \cite{kilian2017}. Satz~\ref{satz:risiko} liefert eine formale
Grundlage f\"ur die These, dass absolute Gew\"ahrleistung f\"ur Software prinzipiell
nicht zugesichert werden kann. Das \textsf{RiskProfile} erg\"anzt dies um eine
diagrammbasierte Nachweisf\"uhrung: Ein vollst\"andig annotiertes UML-Modell, das
$\rho^{\mathrm{sys}}(\mathcal{D}) > 0$ formal dokumentiert, k\"onnte als
prozessualer Nachweis im Sinne einer \emph{Due-Diligence-Verpflichtung} des
Auftragnehmers dienen. Dies hat direkte Implikationen f\"ur:
\begin{itemize}
    \item die Formulierung von Gew\"ahrleistungs- und Haftungsklauseln in AGB und
          Werkvertr\"agen nach \S\S~631~ff.\ BGB,
    \item die Auslegung der EU-Produkthaftungsrichtlinie im Kontext von Softwaresystemen,
    \item die Anforderungen an Risikobewertungen im Rahmen des EU-AI-Acts \cite{euai2024}
          f\"ur KI-basierte Softwaresysteme, wo das \textsf{RiskProfile} als Bestandteil
          der technischen Dokumentation nach Anhang~IV des AI~Acts dienen k\"onnte.
\end{itemize}

% ---------------------------------------------------------------------------
\subsection{Einbettung in Normen und Standards}
% ---------------------------------------------------------------------------

Bestehende Normen wie ISO/IEC~25010 (Softwarequalit\"atsmodell) \cite{iso25010},
ISO/IEC~12207 (Softwarelebenszyklus) und IEC~61508 (funktionale Sicherheit) adressieren
Risiken auf prozessualer Ebene, ohne das strukturelle Restrisiko formal zu begr\"unden.
Eine Erweiterung dieser Normen um ein quantitatives Belastbarkeitsmodell w\"are
w\"unschenswert. Insbesondere k\"onnte die Vertragsklassen-Taxonomie aus
Definition~\ref{def:klassen} als Grundlage f\"ur eine neue Norm zur \emph{Risikokommunikation
in Softwarevertr\"agen} dienen. Das \textsf{RiskProfile} aus Definition~\ref{def:riskprofile}
w\"are dabei ein nat\"urlicher Kandidat f\"ur einen normativen Anhang, der vorschreibt,
welche UML-Diagramme f\"ur welche Vertragsklasse mit Risikoannotationen versehen sein m\"ussen.

% ---------------------------------------------------------------------------
\subsection{Maschinelles Lernen zur Belastbarkeitsvorhersage}
% ---------------------------------------------------------------------------

Moderne KI-Methoden bieten die M\"oglichkeit, den Belastbarkeitsindex $B(t, u, r)$
datengetrieben zu sch\"atzen. Ans\"atze wie \emph{Survival Analysis} (z.\,B.\
Cox-Proportional-Hazards-Modelle \cite{cox1972}) oder \emph{Recurrent Neural Networks}
k\"onnen auf historische Projektdaten trainiert werden, um fr\"uhzeitig Warnsignale
f\"ur einen drohenden Belastbarkeitsverfall zu erkennen. Die in Abschnitt~\ref{sec:uml}
eingef\"uhrten diagrammtypspezifischen Risikoma\ss{}zahlen ---
$\rho^{\mathrm{struct}}$, $\rho^{\mathrm{seq}}$, $\rho^{\mathrm{prop}}$,
$\rho^{\mathrm{state}}$ und $\rho^{\mathrm{ctrl}}$ --- stellen dabei einen
reichhaltigen Feature-Vektor bereit, der als Eingabe f\"ur maschinelle Lernmodelle
dienen kann. Ein solches Modell w\"urde ein
\emph{pr\"adiktives Vertragsmonitoring-System} erm\"oglichen, das Projektmanager
auf Basis des UML-Modells rechtzeitig auf kritische Risikoschwellen hinweist.

% ---------------------------------------------------------------------------
\subsection{Psychologie der Risikowahrnehmung}
% ---------------------------------------------------------------------------

Satz~\ref{satz:risiko} liefert die formale Begr\"undung, dass Softwarerisiken normal
sind. Dennoch zeigt die Praxis, dass Risikooffenbarung in Projekten h\"aufig unterbleibt
--- sei es aus Reputationsangst, kommerziellem Druck oder kognitiven Verzerrungen wie
dem \emph{Optimism Bias} \cite{weinstein1980}. Das \textsf{RiskProfile} k\"onnte hier
als Objektivierungswerkzeug wirken: Wenn die Risikoma\ss{}zahlen automatisch aus dem
Modell berechnet werden (vgl.\ den Ausblick zur automatischen Risikoberechnung), sind
sie nicht mehr das Ergebnis einer subjektiven Einsch\"atzung, sondern folgen formal aus
der Modellstruktur. Zuk\"unftige Arbeiten k\"onnten untersuchen, ob eine solche
Objektivierung die Risikoberichterstattung in Teams verbessert und wie Prozesse gestaltet
werden k\"onnen, die eine offene Risikokultur f\"ordern.

% ---------------------------------------------------------------------------
\subsection{Einfluss k\"unstlicher Intelligenz auf Vertragsbelastbarkeit}
% ---------------------------------------------------------------------------

KI-gest\"utzte Softwareentwicklung (z.\,B.\ durch Large Language Models) ver\"andert die
Dynamik von $u_{\mathrm{tech}}(t)$ erheblich: Einerseits k\"onnen KI-Tools die technische
Unsicherheit in der Entwicklungsphase reduzieren, andererseits f\"uhren sie neue, schwer
verifikable Unsicherheitsquellen ein (Halluzinationen, nicht-deterministische Ausgaben,
mangelnde Erkl\"arbarkeit). F\"ur das \textsf{RiskProfile} ergibt sich daraus eine
besondere Herausforderung: KI-generierter Code besitzt keinen expliziten UML-Entwurf,
aus dem die Risikoma\ss{}zahlen extrahiert werden k\"onnten. Es ist daher eine offene
Forschungsfrage, wie das \textsf{RiskProfile} auf KI-generierte Codebasis angewendet
werden kann --- etwa durch \emph{reverse-engineerte} UML-Modelle aus KI-Code --- und
welche neuen Stereotypen (z.\,B.\ \texttt{\guillemotleft{}riskAIComponent\guillemotright{}})
f\"ur KI-Komponenten in das Profil aufgenommen werden sollten.

% ---------------------------------------------------------------------------
\subsection{Internationale und interdisziplin\"are Vertragsstandards}
% ---------------------------------------------------------------------------

Die unterschiedlichen Rechtssysteme weltweit (Common Law vs.\ kontinentales Recht)
behandeln Softwarevertr\"age und Produkthaftung verschieden. Eine vergleichende Analyse
--- erg\"anzt durch das formale Modell dieser Arbeit und das \textsf{RiskProfile} ---
k\"onnte den Weg f\"ur internationale Mindeststandards zur Risikooffenbarung in
Softwarevertr\"agen ebnen. Das \textsf{RiskProfile} bietet dabei den Vorteil, dass es
sprachunabh\"angig ist: Die formalen Risikoma\ss{}zahlen $\omega(e)$ und das
systemweite Risiko $\rho^{\mathrm{sys}}(\mathcal{D})$ sind unabh\"angig vom
Rechtssystem interpretierbar und k\"onnten daher als gemeinsame technische Grundlage
f\"ur internationale Vertragsstandards dienen, \"ahnlich wie es im Bereich der
Finanzinstrumente durch MiFID~II oder im Medizinbereich durch MDR~2017/745 bereits
etabliert ist.

% ===========================================================================
\begin{thebibliography}{99}

\bibitem{turing1936}
    Turing, A.~M. (1936).
    On Computable Numbers, with an Application to the Entscheidungsproblem.
    \textit{Proceedings of the London Mathematical Society}, 2(42), 230--265.

\bibitem{dijkstra1972}
    Dijkstra, E.~W. (1972).
    The Humble Programmer.
    \textit{Communications of the ACM}, 15(10), 859--866.

\bibitem{brooks1975}
    Brooks, F.~P. (1975).
    \textit{The Mythical Man-Month: Essays on Softwareentwicklering}.
    Addison-Wesley.

\bibitem{hoare1969}
    Hoare, C.~A.~R. (1969).
    An Axiomatic Basis for Computer Programming.
    \textit{Communications of the ACM}, 12(10), 576--580.

\bibitem{cousot1977}
    Cousot, P. \& Cousot, R. (1977).
    Abstract Interpretation: A Unified Lattice Model for Static Analysis of Programs.
    \textit{Proceedings of the 4th ACM POPL}, 238--252.

\bibitem{clarke1999}
    Clarke, E.~M., Grumberg, O., \& Peled, D.~A. (1999).
    \textit{Model Checking}.
    MIT Press.

\bibitem{lamport2002}
    Lamport, L. (2002).
    \textit{Specifying Systems: The TLA+ Language and Tools for Hardware and Softwareentwicklers}.
    Addison-Wesley.

\bibitem{harsanyi1967}
    Harsanyi, J.~C. (1967).
    Games with Incomplete Information Played by Bayesian Players.
    \textit{Management Science}, 14(3), 159--182.

\bibitem{ieee2014}
    IEEE Computer Society (2014).
    \textit{Guide to the Softwareentwicklering Body of Knowledge (SWEBOK), Version 3.0}.
    IEEE Press.

\bibitem{standish2020}
    Standish Group (2020).
    \textit{CHAOS Report 2020: Beyond Infinity}.
    Standish Group International.

\bibitem{isbsg2020}
    ISBSG (2020).
    \textit{ISBSG Benchmark Repository}.
    International Software Benchmarking Standards Group.

\bibitem{iso25010}
    ISO/IEC 25010:2011.
    \textit{Systems and Softwareentwicklering -- Systems and Software Quality Requirements
    and Evaluation (SQuaRE) -- System and Software Quality Models}.
    International Organization for Standardization.

\bibitem{kilian2017}
    Kilian, W. (2017).
    \textit{Computerrechtshandbuch: Informationstechnologie in der Rechts- und
    Wirtschaftspraxis}, 35.~Erg.-Lfg.
    C.~H.~Beck.

\bibitem{euai2024}
    Europ\"aisches Parlament und Rat (2024).
    \textit{Verordnung (EU) 2024/1689 des Europ\"aischen Parlaments und des Rates
    \"uber k\"unstliche Intelligenz (AI Act)}.
    Amtsblatt der Europ\"aischen Union.

\bibitem{cox1972}
    Cox, D.~R. (1972).
    Regression Models and Life-Tables.
    \textit{Journal of the Royal Statistical Society: Series B}, 34(2), 187--202.

\bibitem{weinstein1980}
    Weinstein, N.~D. (1980).
    Unrealistic Optimism About Future Life Events.
    \textit{Journal of Personality and Social Psychology}, 39(5), 806--820.

\bibitem{omg2017}
    Object Management Group (2017).
    \textit{Unified Modeling Language Specification, Version~2.5.1}.
    OMG Document Number: formal/2017-12-05.
    \url{https://www.omg.org/spec/UML/2.5.1}

\bibitem{mccabe1976}
    McCabe, T.~J. (1976).
    A Complexity Measure.
    \textit{IEEE Transactions on Software Engineering}, 2(4), 308--320.

\bibitem{parnas1972}
    Parnas, D.~L. (1972).
    On the Criteria to Be Used in Decomposing Systems into Modules.
    \textit{Communications of the ACM}, 15(12), 1053--1058.

\bibitem{sommerville2011}
    Sommerville, I. (2011).
    \textit{Softwareentwicklering}, 9th~ed.
    Addison-Wesley.

\bibitem{meyer1988}
    Meyer, B. (1988).
    \textit{Object-Oriented Software Construction}.
    Prentice Hall.

\end{thebibliography}

\end{document}
